We aim to characterize involutions and inverses in a more general setting using the microcosm principle. In a certain sense, the data of a \textit{single} mutually invertible pair of arrows or a \textit{single} involution lies within the scope of this principle. This allows one to define relative inverses and involutions which are quite essential in describing dualities of monoidal structures. The first part primarily focus on structures internal to bicategories, whereas the latter part shall explore structures internal to tricategories.

\subsection{\tpdf{$1$-dimensional variants}} \label{sub:1diminv}

We begin by briefly reviewing $0$-dimensional involutions for motivation. As this theory is well known, we omit most of the details. The interested reader should refer to \cite{yauinvcat2020}.

\begin{definition} \leavevmode
	\begin{enumerate}
		\item  A strict involutive object in $\cC$ is given by a morphism $f\colon x \rightarrow x$ such that $f \circ f = 1$.
		\item A morphism of involutions $(x,f) \rightarrow (y,g)$ is an arrow $r\colon x \rightarrow y$ such that $gr=rf$ holds.
	\end{enumerate}
\end{definition}

Involutive objects and their morphisms define a category of involutions $\nInv{0}(\cC)$. This construction extends to a $2$-functor $\nInv{0} \colon \nCat{1} \rightarrow \nCat{1}$, which comes equipped with two canonical natural $2$-transformation, $ U \colon \nInv{0} \Rightarrow 1_{\nCat{1}}$ given by the forgetful action, and its section $i \colon 1_{\nCat{1}} \Rightarrow \nInv{0}$ given by inclusion as identity involution.

Furthermore, the category $\nInv{0}(\nInv{0}(\cC))$ has as objects a pair of involutions $f,g \colon  x \rightarrow x$ such that $fg=gf$. This implies that we have a natural $2$-transformation $\nInv{0}^2 \Rightarrow \nInv{0}$ given by $(f,g) \mapsto (fg)$. In particular, $\nInv{0}$ is a strict $2$-Monad on $\nCat{1}$. 


\begin{example}[$\op_{n,k}$ duality of $\nCat{n}$]
	The $1$-functors $\op_{n,k} \colon \nCat{n}_{\leq 1} \rightarrow \nCat{n}_{\leq 1}$ satisfy the desired axiom $\op_{n,k}^2=1_{\nCat{n}}$. This is a strict $1$-involutive object.
\end{example}

The phrase strict involutive object already suggests the general definition of a weak ($1$-dimensional) involutive object. In particular, we can define weak (or pseudo-) involutive objects inside an arbitrary bicategory $K$. All the corresponding data and axioms shall be written using pasting diagrams, which are to be interpreted via the evaluation functor $e \colon \st K \rightarrow K$. See \cref{appsection:strictification} for more details.


\begin{definition}[pseudo involutive object]
	Let $K$ be a bicategory. We define a weak/pseudo involutive object by the data of
	\begin{enumerate}
		\item a $1$-cell $f \colon x \rightarrow x \in K$,
		\item an invertible $2$-cell
		% https://q.uiver.app/#q=WzAsMyxbMCwxLCJ4Il0sWzEsMCwieCJdLFsyLDEsIngiXSxbMCwxLCJmIl0sWzEsMiwiZiJdLFswLDIsIlxcdmFyZXBzaWxvbl94IiwyXSxbMSw1LCJcXGFscGhhIiwwLHsic2hvcnRlbiI6eyJ0YXJnZXQiOjIwfX1dXQ==
		\[\begin{tikzcd}[ampersand replacement=\&]
			\& x \& \\
			x \&\& x
			\arrow["f", from=1-2, to=2-3]
			\arrow["f", from=2-1, to=1-2]
			\arrow[""{name=0, anchor=center, inner sep=0}, "{\varepsilon_x}"', from=2-1, to=2-3]
			\arrow["\alpha", between={0}{0.8}, Rightarrow, nfold, from=1-2, to=0]
		\end{tikzcd}\cdpunct[8pt]{,}\]
		\item such that the $2$-cell $\alpha$ satisfy an additional axiom as depicted by the following diagram:
		% https://q.uiver.app/#q=WzAsOSxbMSwyLCJ4Il0sWzIsMCwieCJdLFszLDIsIngiXSxbMCwwLCJ4Il0sWzUsMCwieCJdLFs2LDIsIngiXSxbNywwLCJ4Il0sWzgsMiwieCJdLFs0LDEsIlxcZXF1aXYiXSxbMSwyLCJmIl0sWzAsMiwiXFx2YXJlcHNpbG9uX3giLDJdLFszLDAsImYiLDJdLFszLDEsIlxcdmFyZXBzaWxvbl94Il0sWzQsNywiZiIsMV0sWzQsNSwiZiIsMl0sWzQsNiwiXFx2YXJlcHNpbG9uX3giXSxbNiw3LCJmIl0sWzUsNywiXFx2YXJlcHNpbG9uX3giLDJdLFswLDEsImYiXSxbMSwxMCwiXFxhbHBoYSIsMCx7InNob3J0ZW4iOnsic291cmNlIjoyMCwidGFyZ2V0IjoyMH19XSxbMTIsMCwiXFxhbHBoYV57LTF9IiwyLHsic2hvcnRlbiI6eyJzb3VyY2UiOjIwLCJ0YXJnZXQiOjIwfX1dLFs2LDEzLCIiLDEseyJzaG9ydGVuIjp7InNvdXJjZSI6MjAsInRhcmdldCI6MzB9fV0sWzEzLDUsIiIsMSx7InNob3J0ZW4iOnsic291cmNlIjoyMCwidGFyZ2V0IjoyMH19XV0=
		\[\begin{tikzcd}[ampersand replacement=\&]
			x \&\& x \&\&\& x \&\& x \& \\
			\&\&\&\& \equiv \\
			\& x \&\& x \&\&\& x \&\& x
			\arrow[""{name=0, anchor=center, inner sep=0}, "{\varepsilon_x}", from=1-1, to=1-3]
			\arrow["f"', from=1-1, to=3-2]
			\arrow["f", from=1-3, to=3-4]
			\arrow["{\varepsilon_x}", from=1-6, to=1-8]
			\arrow["f"', from=1-6, to=3-7]
			\arrow[""{name=1, anchor=center, inner sep=0}, "f"{description}, from=1-6, to=3-9]
			\arrow["f", from=1-8, to=3-9]
			\arrow["f", from=3-2, to=1-3]
			\arrow[""{name=2, anchor=center, inner sep=0}, "{\varepsilon_x}"', from=3-2, to=3-4]
			\arrow["{\varepsilon_x}"', from=3-7, to=3-9]
			\arrow["{\alpha^{-1}}"', between={0.2}{0.8}, Rightarrow, nfold, from=0, to=3-2]
			\arrow["\alpha", between={0.2}{0.8}, Rightarrow, nfold, from=1-3, to=2]
			\arrow[between={0.2}{0.8}, Rightarrow, nfold, from=1, to=3-7]
			\arrow[between={0.2}{0.7}, Rightarrow, nfold, from=1-8, to=1]
		\end{tikzcd}\cdpunct[8pt]{.}\]
	\end{enumerate}
\end{definition}

\begin{example}
	Every strict involutive object (where $f^2=1$) in a strict $2$-category is trivially a weak involutive object.
\end{example}

One might wonder why do we require the additional identities in the specification of a weak involutive object. The key point is that the structure is incoherent without the additional axiom, analogous to how specifying a monoidal category without its pentagon or triangle identity is meaningless. This extra identity ensures that all reasonable ways of going from $f^n \Rightarrow f^m$ are equivalent. In fact, there are four equivalent way of expressing this additional axiom. The proof is relatively simple and uses the pasting of the $2$-cells $\alpha,\alpha^{-1}$, or the unitors appropriately.

\begin{lemma}
	The data of an involution in a bicategory is equivalently expressed as a $1$-cell $f\colon x \rightarrow x$ and an invertible $2$-cell $\alpha \colon f^2 \Rightarrow 1$ satisfying any one of the following axioms:
	\begin{enumerate}
		\item % https://q.uiver.app/#q=WzAsOSxbMSwyLCJ4Il0sWzIsMCwieCJdLFszLDIsIngiXSxbMCwwLCJ4Il0sWzUsMCwieCJdLFs2LDIsIngiXSxbNywwLCJ4Il0sWzgsMiwieCJdLFs0LDEsIlxcZXF1aXYiXSxbMSwyLCJmIl0sWzAsMiwiXFx2YXJlcHNpbG9uX3giLDJdLFszLDAsImYiLDJdLFszLDEsIlxcdmFyZXBzaWxvbl94Il0sWzQsNywiZiIsMV0sWzQsNSwiZiIsMl0sWzQsNiwiXFx2YXJlcHNpbG9uX3giXSxbNiw3LCJmIl0sWzUsNywiXFx2YXJlcHNpbG9uX3giLDJdLFswLDEsImYiXSxbMSwxMCwiXFxhbHBoYSIsMCx7InNob3J0ZW4iOnsic291cmNlIjoyMCwidGFyZ2V0IjoyMH19XSxbMTIsMCwiXFxhbHBoYV57LTF9IiwyLHsic2hvcnRlbiI6eyJzb3VyY2UiOjIwLCJ0YXJnZXQiOjIwfX1dLFs2LDEzLCIiLDEseyJzaG9ydGVuIjp7InNvdXJjZSI6MjAsInRhcmdldCI6MzB9fV0sWzEzLDUsIiIsMSx7InNob3J0ZW4iOnsic291cmNlIjoyMCwidGFyZ2V0IjoyMH19XV0=
		\[\begin{tikzcd}[ampersand replacement=\&]
			x \&\& x \&\&\& x \&\& x \& \\
			\&\&\&\& \equiv \\
			\& x \&\& x \&\&\& x \&\& x
			\arrow[""{name=0, anchor=center, inner sep=0}, "{\varepsilon_x}", from=1-1, to=1-3]
			\arrow["f"', from=1-1, to=3-2]
			\arrow["f", from=1-3, to=3-4]
			\arrow["{\varepsilon_x}", from=1-6, to=1-8]
			\arrow["f"', from=1-6, to=3-7]
			\arrow[""{name=1, anchor=center, inner sep=0}, "f"{description}, from=1-6, to=3-9]
			\arrow["f", from=1-8, to=3-9]
			\arrow["f", from=3-2, to=1-3]
			\arrow[""{name=2, anchor=center, inner sep=0}, "{\varepsilon_x}"', from=3-2, to=3-4]
			\arrow["{\varepsilon_x}"', from=3-7, to=3-9]
			\arrow["{\alpha^{-1}}"', between={0.2}{0.8}, Rightarrow, nfold, from=0, to=3-2]
			\arrow["\alpha", between={0.2}{0.8}, Rightarrow, nfold, from=1-3, to=2]
			\arrow[between={0.2}{0.8}, Rightarrow, nfold, from=1, to=3-7]
			\arrow[between={0.2}{0.7}, Rightarrow, nfold, from=1-8, to=1]
		\end{tikzcd}\cdpunct[8pt]{.}\]
		\item % https://q.uiver.app/#q=WzAsOSxbMSwxLCJ4Il0sWzMsMCwieCJdLFs0LDIsIngiXSxbMCwyLCJ4Il0sWzcsMCwieCJdLFsxMCwyLCJ4Il0sWzksMSwieCJdLFs2LDIsIngiXSxbNSwxLCJcXGVxdWl2Il0sWzAsMSwiZiJdLFsxLDIsImYiXSxbMCwyLCJcXHZhcmVwc2lsb25feCIsMV0sWzMsMCwiZiJdLFs2LDUsImYiXSxbNCw2LCJmIl0sWzcsNCwiZiJdLFs3LDYsIlxcdmFyZXBzaWxvbl94IiwxXSxbMywyLCJmIiwyXSxbNyw1LCJmIiwyXSxbNCwxNiwiXFxhbHBoYSIsMCx7InNob3J0ZW4iOnsidGFyZ2V0IjoyMH19XSxbMCwxNywiIiwwLHsic2hvcnRlbiI6eyJzb3VyY2UiOjIwLCJ0YXJnZXQiOjIwfX1dLFs2LDE4LCIiLDAseyJzaG9ydGVuIjp7InNvdXJjZSI6MjAsInRhcmdldCI6MjB9fV0sWzEsMTEsIlxcYWxwaGEiLDAseyJzaG9ydGVuIjp7InRhcmdldCI6MjB9fV1d
		\[\begin{tikzcd}[ampersand replacement=\&]
			\&\&\& x \&\&\&\& x \&\&\& \\
			\& x \&\&\&\& \equiv \&\&\&\& x \\
			x \&\&\&\& x \&\& x \&\&\&\& x
			\arrow["f", from=1-4, to=3-5]
			\arrow["f", from=1-8, to=2-10]
			\arrow["f", from=2-2, to=1-4]
			\arrow[""{name=0, anchor=center, inner sep=0}, "{\varepsilon_x}"{description}, from=2-2, to=3-5]
			\arrow["f", from=2-10, to=3-11]
			\arrow["f", from=3-1, to=2-2]
			\arrow[""{name=1, anchor=center, inner sep=0}, "f"', from=3-1, to=3-5]
			\arrow["f", from=3-7, to=1-8]
			\arrow[""{name=2, anchor=center, inner sep=0}, "{\varepsilon_x}"{description}, from=3-7, to=2-10]
			\arrow[""{name=3, anchor=center, inner sep=0}, "f"', from=3-7, to=3-11]
			\arrow["\alpha", between={0}{0.8}, Rightarrow, nfold, from=1-4, to=0]
			\arrow["\alpha", between={0}{0.8}, Rightarrow, nfold, from=1-8, to=2]
			\arrow[between={0.2}{0.8}, Rightarrow, nfold, from=2-2, to=1]
			\arrow[between={0.2}{0.8}, Rightarrow, nfold, from=2-10, to=3]
		\end{tikzcd}\cdpunct[8pt]{.}\]
		\item % https://q.uiver.app/#q=WzAsOSxbMCwyLCJ4Il0sWzEsMCwieCJdLFsyLDIsIngiXSxbMywwLCJ4Il0sWzUsMiwieCJdLFs2LDAsIngiXSxbNywyLCJ4Il0sWzgsMCwieCJdLFs0LDEsIlxcZXF1aXYiXSxbMCwxLCJmIl0sWzIsMywiZiIsMl0sWzEsMiwiZiJdLFswLDIsIlxcdmFyZXBzaWxvbl94IiwyXSxbMSwzLCJcXHZhcmVwc2lsb25feCJdLFs0LDcsImYiLDFdLFs0LDUsImYiXSxbNSw3LCJcXHZhcmVwc2lsb25feCJdLFs2LDcsImYiLDJdLFs0LDYsIlxcdmFyZXBzaWxvbl94IiwyXSxbMSwxMiwiXFxhbHBoYSIsMCx7InNob3J0ZW4iOnsic291cmNlIjoyMCwidGFyZ2V0IjoyMH19XSxbNSwxNCwiIiwxLHsic2hvcnRlbiI6eyJzb3VyY2UiOjIwLCJ0YXJnZXQiOjQwfX1dLFsxNCw2LCIiLDEseyJzaG9ydGVuIjp7InNvdXJjZSI6NDAsInRhcmdldCI6MjB9fV0sWzEzLDIsIlxcYWxwaGFeey0xfSIsMCx7InNob3J0ZW4iOnsic291cmNlIjoyMCwidGFyZ2V0IjoyMH19XV0=
		\[\begin{tikzcd}[ampersand replacement=\&]
			\& x \&\& x \&\&\& x \&\& x \\
			\&\&\&\& \equiv \\
			x \&\& x \&\&\& x \&\& x
			\arrow[""{name=0, anchor=center, inner sep=0}, "{\varepsilon_x}", from=1-2, to=1-4]
			\arrow["f", from=1-2, to=3-3]
			\arrow["{\varepsilon_x}", from=1-7, to=1-9]
			\arrow["f", from=3-1, to=1-2]
			\arrow[""{name=1, anchor=center, inner sep=0}, "{\varepsilon_x}"', from=3-1, to=3-3]
			\arrow["f"', from=3-3, to=1-4]
			\arrow["f", from=3-6, to=1-7]
			\arrow[""{name=2, anchor=center, inner sep=0}, "f"{description}, from=3-6, to=1-9]
			\arrow["{\varepsilon_x}"', from=3-6, to=3-8]
			\arrow["f"', from=3-8, to=1-9]
			\arrow["\alpha", between={0.2}{0.8}, Rightarrow, nfold, from=1-2, to=1]
			\arrow["{\alpha^{-1}}", between={0.2}{0.8}, Rightarrow, nfold, from=0, to=3-3]
			\arrow[between={0.2}{0.6}, Rightarrow, nfold, from=1-7, to=2]
			\arrow[between={0.4}{0.8}, Rightarrow, nfold, from=2, to=3-8]
		\end{tikzcd}\cdpunct[8pt]{.}\]
		\item % https://q.uiver.app/#q=WzAsOSxbMSwxLCJ4Il0sWzMsMywieCJdLFs0LDEsIngiXSxbMCwwLCJ4Il0sWzcsMywieCJdLFsxMCwwLCJ4Il0sWzksMSwieCJdLFs2LDEsIngiXSxbNSwyLCJcXGVxdWl2Il0sWzAsMSwiZiIsMl0sWzEsMiwiZiIsMl0sWzAsMiwiXFx2YXJlcHNpbG9uX3giLDFdLFszLDAsImYiLDJdLFs2LDUsImYiLDJdLFs0LDYsImYiLDJdLFs3LDQsImYiLDJdLFs3LDYsIlxcdmFyZXBzaWxvbl94IiwxXSxbMywyLCJmIl0sWzcsNSwiZiJdLFsxMSwxLCJcXGFscGhhXnstMX0iLDIseyJzaG9ydGVuIjp7InNvdXJjZSI6MjAsInRhcmdldCI6MjB9fV0sWzE2LDQsIlxcYWxwaGFeey0xfSIsMix7InNob3J0ZW4iOnsic291cmNlIjoyMCwidGFyZ2V0IjoyMH19XSxbMTcsMCwiIiwyLHsic2hvcnRlbiI6eyJzb3VyY2UiOjIwLCJ0YXJnZXQiOjIwfX1dLFsxOCw2LCIiLDIseyJzaG9ydGVuIjp7InNvdXJjZSI6MjAsInRhcmdldCI6MjB9fV1d
		\[\begin{tikzcd}[ampersand replacement=\&]
			x \&\&\&\&\&\&\&\&\&\& x \\
			\& x \&\&\& x \&\& x \&\&\& x \\
			\&\&\&\&\& \equiv \\
			\&\&\& x \&\&\&\& x
			\arrow["f"', from=1-1, to=2-2]
			\arrow[""{name=0, anchor=center, inner sep=0}, "f", from=1-1, to=2-5]
			\arrow[""{name=1, anchor=center, inner sep=0}, "{\varepsilon_x}"{description}, from=2-2, to=2-5]
			\arrow["f"', from=2-2, to=4-4]
			\arrow[""{name=2, anchor=center, inner sep=0}, "f", from=2-7, to=1-11]
			\arrow[""{name=3, anchor=center, inner sep=0}, "{\varepsilon_x}"{description}, from=2-7, to=2-10]
			\arrow["f"', from=2-7, to=4-8]
			\arrow["f"', from=2-10, to=1-11]
			\arrow["f"', from=4-4, to=2-5]
			\arrow["f"', from=4-8, to=2-10]
			\arrow[between={0.2}{0.8}, Rightarrow, nfold, from=0, to=2-2]
			\arrow["{\alpha^{-1}}"', between={0.2}{0.8}, Rightarrow, nfold, from=1, to=4-4]
			\arrow[between={0.2}{0.8}, Rightarrow, nfold, from=2, to=2-10]
			\arrow["{\alpha^{-1}}"', between={0.2}{0.8}, Rightarrow, nfold, from=3, to=4-8]
		\end{tikzcd}\cdpunct[8pt]{.}\]
	\end{enumerate} 
\end{lemma}

\begin{comment}
	\begin{proof}
		All of these diagrams are derivable by using appropriate pasting with the $2$-cells $\alpha$, $\alpha^{-1}$ and the unitors. For instance, the involution definition implying the axiom $1$ can be observed from the following commuting diagram:
		% https://q.uiver.app/#q=WzAsMjgsWzEsNCwieCJdLFsyLDIsIngiXSxbMyw0LCJ4Il0sWzAsMiwieCJdLFs1LDIsIngiXSxbNiw0LCJ4Il0sWzcsMiwieCJdLFs4LDQsIngiXSxbNCwzLCJcXGVxdWl2Il0sWzEsMCwieCJdLFs2LDAsIngiXSxbMCw5LCJ4Il0sWzEsOCwieCJdLFszLDgsIngiXSxbMiw3LCJ4Il0sWzIsNSwiXFx2ZXF1aXYiXSxbNyw1LCJcXHZlcXVpdiJdLFs0LDcsIlxcZXF1aXYiXSxbNiw4LCJ4Il0sWzgsMTAsIngiXSxbNyw5LCJ4Il0sWzgsNywieCJdLFs3LDYsIngiXSxbNCw5XSxbNCwxMF0sWzQsMF0sWzQsMV0sWzQsMTEsIlxcY2RvdHMiXSxbMSwyLCJmIl0sWzAsMiwiMV94IiwyXSxbMywwLCJmIiwyXSxbMywxLCIxX3giXSxbNCw3LCJmIiwxXSxbNCw1LCJmIiwyXSxbNCw2LCIxX3giXSxbNiw3LCJmIl0sWzUsNywiMV94IiwyXSxbMCwxLCJmIl0sWzMsOSwiZiJdLFs5LDEsImYiXSxbNCwxMCwiZiJdLFsxMCw2LCJmIl0sWzExLDEyLCJmIiwyXSxbMTIsMTQsImYiXSxbMTIsMTMsIjFfeCIsMl0sWzE0LDEzLCJmIl0sWzE4LDE5LCJmIiwxLHsib2Zmc2V0IjozLCJjdXJ2ZSI6LTUsInNob3J0ZW4iOnsic291cmNlIjoxMCwidGFyZ2V0IjoxMH19XSxbMTgsMjAsImYiLDJdLFsyMCwxOSwiMV94IiwyXSxbMTgsMjEsIjFfeCIsMl0sWzIxLDE5LCJmIl0sWzE4LDIyLCJmIl0sWzIyLDIxLCJmIl0sWzIzLDI0LCJcXHRleHRybXtwYXN0ZSB1bml0b3JzfSIsMCx7ImxldmVsIjozfV0sWzI1LDI2LCJcXHRleHRybXtwYXN0ZSB9XFxhbHBoYSIsMCx7ImxldmVsIjozfV0sWzEsMjksIlxcYWxwaGEiLDAseyJzaG9ydGVuIjp7InNvdXJjZSI6MjAsInRhcmdldCI6MjB9fV0sWzMxLDAsIlxcYWxwaGFeey0xfSIsMix7InNob3J0ZW4iOnsic291cmNlIjoyMCwidGFyZ2V0IjoyMH19XSxbNiwzMiwiIiwxLHsic2hvcnRlbiI6eyJzb3VyY2UiOjIwLCJ0YXJnZXQiOjMwfX1dLFszMiw1LCIiLDEseyJzaG9ydGVuIjp7InNvdXJjZSI6MjAsInRhcmdldCI6MjB9fV0sWzksMzEsIlxcYWxwaGEiLDIseyJzaG9ydGVuIjp7InNvdXJjZSI6MjAsInRhcmdldCI6MjB9fV0sWzEwLDM0LCJcXGFscGhhIiwwLHsic2hvcnRlbiI6eyJzb3VyY2UiOjIwLCJ0YXJnZXQiOjIwfX1dLFsyMiw0OSwiXFxhbHBoYSIsMix7InNob3J0ZW4iOnsic291cmNlIjoyMCwidGFyZ2V0IjoyMH19XSxbMjEsNDYsIiIsMSx7InNob3J0ZW4iOnsic291cmNlIjoyMCwidGFyZ2V0IjoyMH19XSxbNDYsMjAsIiIsMSx7InNob3J0ZW4iOnsic291cmNlIjo1MH19XSxbMTQsNDQsIlxcYWxwaGEiLDIseyJzaG9ydGVuIjp7InNvdXJjZSI6MjAsInRhcmdldCI6MjB9fV1d
		\[\begin{tikzcd}[ampersand replacement=\&]
			\& x \&\&\& {} \&\& x \&\& \\
			\&\&\&\& {} \\
			x \&\& x \&\&\& x \&\& x \\
			\&\&\&\& \equiv \\
			\& x \&\& x \&\&\& x \&\& x \\
			\&\& \vequiv \&\&\&\&\& \vequiv \\
			\&\&\&\&\&\&\& x \\
			\&\& x \&\& \equiv \&\&\&\& x \\
			\& x \&\& x \&\&\& x \\
			x \&\&\&\& {} \&\&\& x \\
			\&\&\&\& {} \&\&\&\& x \\
			\&\&\&\& \cdots
			\arrow["f", from=1-2, to=3-3]
			\arrow["{\textrm{paste }\alpha}", Rightarrow, scaling nfold=3, from=1-5, to=2-5]
			\arrow["f", from=1-7, to=3-8]
			\arrow["f", from=3-1, to=1-2]
			\arrow[""{name=0, anchor=center, inner sep=0}, "{1_x}", from=3-1, to=3-3]
			\arrow["f"', from=3-1, to=5-2]
			\arrow["f", from=3-3, to=5-4]
			\arrow["f", from=3-6, to=1-7]
			\arrow[""{name=1, anchor=center, inner sep=0}, "{1_x}", from=3-6, to=3-8]
			\arrow["f"', from=3-6, to=5-7]
			\arrow[""{name=2, anchor=center, inner sep=0}, "f"{description}, from=3-6, to=5-9]
			\arrow["f", from=3-8, to=5-9]
			\arrow["f", from=5-2, to=3-3]
			\arrow[""{name=3, anchor=center, inner sep=0}, "{1_x}"', from=5-2, to=5-4]
			\arrow["{1_x}"', from=5-7, to=5-9]
			\arrow["f", from=7-8, to=8-9]
			\arrow["f", from=8-3, to=9-4]
			\arrow["f", from=8-9, to=11-9]
			\arrow["f", from=9-2, to=8-3]
			\arrow[""{name=4, anchor=center, inner sep=0}, "{1_x}"', from=9-2, to=9-4]
			\arrow["f", from=9-7, to=7-8]
			\arrow[""{name=5, anchor=center, inner sep=0}, "{1_x}"', from=9-7, to=8-9]
			\arrow["f"', from=9-7, to=10-8]
			\arrow[""{name=6, anchor=center, inner sep=0}, "f"{description}, shift right=3, curve={height=-30pt}, between={0.1}{0.9}, from=9-7, to=11-9]
			\arrow["f"', from=10-1, to=9-2]
			\arrow["{\textrm{paste unitors}}", Rightarrow, scaling nfold=3, from=10-5, to=11-5]
			\arrow["{1_x}"', from=10-8, to=11-9]
			\arrow["\alpha"', between={0.2}{0.8}, Rightarrow, nfold, from=1-2, to=0]
			\arrow["\alpha", between={0.2}{0.8}, Rightarrow, nfold, from=1-7, to=1]
			\arrow["{\alpha^{-1}}"', between={0.2}{0.8}, Rightarrow, nfold, from=0, to=5-2]
			\arrow["\alpha", between={0.2}{0.8}, Rightarrow, nfold, from=3-3, to=3]
			\arrow[between={0.2}{0.8}, Rightarrow, nfold, from=2, to=5-7]
			\arrow[between={0.2}{0.7}, Rightarrow, nfold, from=3-8, to=2]
			\arrow["\alpha"', between={0.2}{0.8}, Rightarrow, nfold, from=7-8, to=5]
			\arrow["\alpha"', between={0.2}{0.8}, Rightarrow, nfold, from=8-3, to=4]
			\arrow[between={0.2}{0.8}, Rightarrow, nfold, from=8-9, to=6]
			\arrow[between={0.5}{1}, Rightarrow, nfold, from=6, to=10-8]
		\end{tikzcd}\cdpunct[8pt]{.}\]
		The rest can be derived similarly. 
	\end{proof}
\end{comment}

The $2$nd and the $4$th form of the additional axiom are particularly interesting as they allow one to define a lax/colax variant for involutive object in a (lax/colax) bicategory. Thus, more generally we can work with \emph{lax involutive objects}. It is now instructive to define $1$-morphisms and $2$-morphisms of these involutions explicitly. 

\begin{definition}
	Let $(x,f,\alpha)$ and $(y,g,\beta)$ be two (lax) involutive objects in a bicategory $K$. A (lax) $1$-morphism of involution $(r,u) \colon (x,f,\alpha) \rightarrow (y,g,\beta)$ is
	\begin{enumerate}
		\item a $1$-cell $r \colon x \rightarrow y \in K$ 
		\item a $2$-cell 
		% https://q.uiver.app/#q=WzAsNCxbMCwwLCJ4Il0sWzIsMCwieSJdLFswLDIsIngiXSxbMiwyLCJ5Il0sWzAsMiwiZiIsMl0sWzEsMywiZyJdLFswLDEsInIiXSxbMiwzLCJyIiwyXSxbMSwyLCJ1IiwwLHsic2hvcnRlbiI6eyJzb3VyY2UiOjIwLCJ0YXJnZXQiOjIwfSwibGV2ZWwiOjJ9XV0=
		\[\begin{tikzcd}[ampersand replacement=\&]
			x \&\& y \\
			\\
			x \&\& y
			\arrow["r", from=1-1, to=1-3]
			\arrow["f"', from=1-1, to=3-1]
			\arrow["u", between={0.2}{0.8}, Rightarrow, nfold, from=1-3, to=3-1]
			\arrow["g", from=1-3, to=3-3]
			\arrow["r"', from=3-1, to=3-3]
		\end{tikzcd}\cdpunct[8pt]{,}\]
		\item such that the following axiom holds:
		% https://q.uiver.app/#q=WzAsMTIsWzAsMCwieCJdLFsyLDAsInkiXSxbMSwyLCJ4Il0sWzMsMiwieSJdLFswLDQsIngiXSxbMiw0LCJ5Il0sWzQsMiwiXFxlcXVpdiJdLFs1LDEsIngiXSxbNSwzLCJ4Il0sWzcsMSwieSJdLFs3LDMsInkiXSxbOSwyLCJ5Il0sWzAsMiwiZiJdLFsxLDMsImciXSxbMCwxLCJyIl0sWzIsMywiciIsMl0sWzEsMiwidSIsMCx7InNob3J0ZW4iOnsic291cmNlIjoyMCwidGFyZ2V0IjoyMH0sImxldmVsIjoyfV0sWzMsNSwiZyJdLFsyLDQsImYiXSxbNCw1LCJyIiwyXSxbMyw0LCJ1IiwwLHsic2hvcnRlbiI6eyJzb3VyY2UiOjIwLCJ0YXJnZXQiOjIwfSwibGV2ZWwiOjJ9XSxbOCwxMCwiciIsMl0sWzcsOCwiXFx2YXJlcHNpbG9uX3giLDJdLFs3LDksInIiXSxbOSwxMCwiXFx2YXJlcHNpbG9uX3kiLDJdLFs5LDExLCJnIl0sWzExLDEwLCJnIl0sWzcsMTAsInIiLDFdLFswLDQsIlxcdmFyZXBzaWxvbl94IiwyXSxbMTEsMjQsIlxcYmV0YSIsMix7InNob3J0ZW4iOnsic291cmNlIjoyMCwidGFyZ2V0IjoyMH19XSxbMjcsOCwiIiwxLHsic2hvcnRlbiI6eyJzb3VyY2UiOjIwLCJ0YXJnZXQiOjIwfX1dLFs5LDI3LCIiLDEseyJzaG9ydGVuIjp7InNvdXJjZSI6MjAsInRhcmdldCI6MjB9fV0sWzIsMjgsIlxcYWxwaGEiLDIseyJzaG9ydGVuIjp7InNvdXJjZSI6MjAsInRhcmdldCI6MjB9fV1d
		\[\begin{tikzcd}[ampersand replacement=\&]
			x \&\& y \&\&\&\&\&\&\& \\
			\&\&\&\&\& x \&\& y \\
			\& x \&\& y \& \equiv \&\&\&\&\& y \\
			\&\&\&\&\& x \&\& y \\
			x \&\& y
			\arrow["r", from=1-1, to=1-3]
			\arrow["f", from=1-1, to=3-2]
			\arrow[""{name=0, anchor=center, inner sep=0}, "{\varepsilon_x}"', from=1-1, to=5-1]
			\arrow["u", between={0.2}{0.8}, Rightarrow, nfold, from=1-3, to=3-2]
			\arrow["g", from=1-3, to=3-4]
			\arrow["r", from=2-6, to=2-8]
			\arrow["{\varepsilon_x}"', from=2-6, to=4-6]
			\arrow[""{name=1, anchor=center, inner sep=0}, "r"{description}, from=2-6, to=4-8]
			\arrow["g", from=2-8, to=3-10]
			\arrow[""{name=2, anchor=center, inner sep=0}, "{\varepsilon_y}"', from=2-8, to=4-8]
			\arrow["r"', from=3-2, to=3-4]
			\arrow["f", from=3-2, to=5-1]
			\arrow["u", between={0.2}{0.8}, Rightarrow, nfold, from=3-4, to=5-1]
			\arrow["g", from=3-4, to=5-3]
			\arrow["g", from=3-10, to=4-8]
			\arrow["r"', from=4-6, to=4-8]
			\arrow["r"', from=5-1, to=5-3]
			\arrow[between={0.2}{0.8}, Rightarrow, nfold, from=1, to=4-6]
			\arrow[between={0.2}{0.8}, Rightarrow, nfold, from=2-8, to=1]
			\arrow["\alpha"', between={0.2}{0.8}, Rightarrow, nfold, from=3-2, to=0]
			\arrow["\beta"', between={0.2}{0.8}, Rightarrow, nfold, from=3-10, to=2]
		\end{tikzcd}\cdpunct[8pt]{.}\]
	\end{enumerate}
\end{definition}

\begin{definition}
	Let $(r,u)$ and $(t,v)$ be two $1$-morphisms of involution between $(x,f,\alpha)$ and $(y,g,\beta)$. A $2$-morphism $\gamma \colon (r,u) \Rightarrow (t,v)$ is
	
	\begin{enumerate}
		\item a $2$-cell $r \Rightarrow t \in K$,
		
		\item which satisfy the following axiom
		% https://q.uiver.app/#q=WzAsOSxbMCwwLCJ4Il0sWzAsMiwieCJdLFsyLDAsInkiXSxbMiwyLCJ5Il0sWzMsMSwiXFxlcXVpdiJdLFs0LDAsIngiXSxbNCwyLCJ4Il0sWzYsMCwieSJdLFs2LDIsInkiXSxbMCwxLCJmIiwyXSxbMSwzLCJyIl0sWzIsMywiZyJdLFswLDIsInIiLDJdLFsyLDEsInUiLDAseyJzaG9ydGVuIjp7InNvdXJjZSI6MjAsInRhcmdldCI6MjB9LCJsZXZlbCI6Mn1dLFs1LDcsInQiLDJdLFs1LDYsImYiLDJdLFs3LDgsImciXSxbNiw4LCJ0Il0sWzcsNiwidiIsMCx7InNob3J0ZW4iOnsic291cmNlIjoyMCwidGFyZ2V0IjoyMH0sImxldmVsIjoyfV0sWzEsMywidCIsMix7ImN1cnZlIjo1fV0sWzUsNywiciIsMCx7ImN1cnZlIjotNX1dLFsxMCwxOSwiXFxnYW1tYSIsMCx7InNob3J0ZW4iOnsic291cmNlIjoyMCwidGFyZ2V0IjoyMH19XSxbMjAsMTQsIlxcZ2FtbWEiLDAseyJzaG9ydGVuIjp7InNvdXJjZSI6MjAsInRhcmdldCI6MjB9fV1d
		\[\begin{tikzcd}[ampersand replacement=\&]
			x \&\& y \&\& x \&\& y \\
			\&\&\& \equiv \\
			x \&\& y \&\& x \&\& y
			\arrow["r"', from=1-1, to=1-3]
			\arrow["f"', from=1-1, to=3-1]
			\arrow["u", between={0.2}{0.8}, Rightarrow, nfold, from=1-3, to=3-1]
			\arrow["g", from=1-3, to=3-3]
			\arrow[""{name=0, anchor=center, inner sep=0}, "t"', from=1-5, to=1-7]
			\arrow[""{name=1, anchor=center, inner sep=0}, "r", curve={height=-30pt}, from=1-5, to=1-7]
			\arrow["f"', from=1-5, to=3-5]
			\arrow["v", between={0.2}{0.8}, Rightarrow, nfold, from=1-7, to=3-5]
			\arrow["g", from=1-7, to=3-7]
			\arrow[""{name=2, anchor=center, inner sep=0}, "r", from=3-1, to=3-3]
			\arrow[""{name=3, anchor=center, inner sep=0}, "t"', curve={height=30pt}, from=3-1, to=3-3]
			\arrow["t", from=3-5, to=3-7]
			\arrow["\gamma", between={0.2}{0.8}, Rightarrow, nfold, from=1, to=0]
			\arrow["\gamma", between={0.2}{0.8}, Rightarrow, nfold, from=2, to=3]
		\end{tikzcd}\cdpunct[8pt]{.}\]
	\end{enumerate}
\end{definition}

As usual, all of this data assembles into a \emph{bicategory of involutions} $\nInv{1}(K)$ associated to any bicategory $K$. This follows by simply defining all the relevant compositional data strictly in $\st K$ and then descending those diagrams onto $K$ using the evaluation functor $e \colon \st K \rightarrow K$ appropriately. We do note that the identity $1$-morphisms in $\nInv{1}(K)$ are given by the $1$-cells $1_x \colon x\rightarrow x$, which are not the empty strings.

\begin{lemma}
	There is a bicategory of involutions $\nInv{1}(K)$ corresponding to each bicategory $K$.
\end{lemma}

\begin{remark}
	With a lengthy computation and using strifications appropriately, one could show that $\nInv{1}$ is a (strict) trifunctor  $\nInv{1} \colon \BiCat \rightarrow \BiCat$ on the tricategory of $\BiCat$. One could also show that there is a (strict) tritransformation $U \colon \nInv{1} \Rightarrow 1_{\BiCat}$ and a (strict) tritransformation $i \colon 1_{\BiCat} \Rightarrow \nInv{1}$. However, we shall not go into the details here.
	%strification should commute past 2inv in reality
\end{remark}

We have now provided enough discussion about the theory of involutions and want to motivate a more general definition of an involutive object. Consider again the $2$-functor $\op_1 \colon \op_2(\nCat{1}) \rightarrow \nCat{1}$. Here, the lower dimensional involution has its domain altered by a higher dimensional involution. This suggests the definition of a relative involutive object. A similar account of these can also be found in \cite{yauinvcat2020}. Hence, we shall be brief in describing the known details and provide a proof only for the missing slice construction.

\begin{definition}
	Let $(\cC,T,\alpha)$ be a pseudo (or lax) involutive object in $\nCat{1}$. We define an involutive object relative to $T$ as a morphism $f \colon Tx \rightarrow x \in \cC$ such that the following diagram commutes:
	% https://q.uiver.app/#q=WzAsNCxbMCwwLCJUXjJ4Il0sWzIsMCwiVHgiXSxbMiwyLCJ4Il0sWzAsMiwieCJdLFsxLDIsImYiXSxbMCwxLCJUZiJdLFswLDMsIlxcYWxwaGFfeCIsMl0sWzMsMiwiMSIsMl1d
	\[\begin{tikzcd}[ampersand replacement=\&]
		{T^2x} \&\& Tx \\
		\\
		x \&\& x
		\arrow["Tf", from=1-1, to=1-3]
		\arrow["{\alpha_x}"', from=1-1, to=3-1]
		\arrow["f", from=1-3, to=3-3]
		\arrow["1"', from=3-1, to=3-3]
	\end{tikzcd}\cdpunct[8pt]{.}\]
\end{definition}

\begin{remark}
	To relate with the previous definition, we define the adjunct morphism by $f^{-1} \colon x \xrightarrow{\alpha^{-1}_x} T^2x \xrightarrow{Tf} Tx$. From this, it follows that both $f^{-1} \circ f = 1 $ and $1 = f \circ f^{-1}$ are satisfied. 
	Furthermore, if one of the equation is satisfied then the other one is automatically satisfied due to the additional axiom of $T$. However, the invertibility of $\alpha$ is essential for this, as $f^{-1}$ is not defined without $\alpha^{-1}$.
\end{remark}

\begin{lemma}
	Let $(C,T,\alpha)$ be a lax involutive object in $\nCat{1}$ and $(x,f \colon Tx \rightarrow x)$ be an involutive object relative to $T$. Then $C/x$ admits a canonical lax involutive structure using:
	% https://q.uiver.app/#q=WzAsMTMsWzAsMCwiQy94Il0sWzIsMCwiQy94Il0sWzAsMSwiYyJdLFswLDIsIngiXSxbMiwxLCJUYyJdLFsyLDIsIlR4Il0sWzIsMywieCJdLFszLDAsIlxcbGVmdChcXGFscGhhX3svfVxccmlnaHQpX3tjfSA6Il0sWzQsMSwiVF4yeCJdLFs0LDIsIlR4Il0sWzQsMCwiVF4yYyJdLFs2LDMsIngiXSxbNiwwLCJjIl0sWzAsMSwiVC9mIl0sWzIsM10sWzQsNV0sWzUsNiwiZiJdLFs4LDExLCJcXGFscGhhX3giXSxbOCw5LCJUZiIsMl0sWzksMTEsImYiLDJdLFsxMiwxMV0sWzEwLDEyLCJcXGFscGhhX2MiXSxbMTAsOF0sWzE0LDE1LCIiLDAseyJzaG9ydGVuIjp7InNvdXJjZSI6MjAsInRhcmdldCI6MjB9LCJsZXZlbCI6MSwic3R5bGUiOnsidGFpbCI6eyJuYW1lIjoibWFwcyB0byJ9fX1dXQ==
	\[\begin{tikzcd}[ampersand replacement=\&]
		{C/x} \&\& {C/x} \& {\left(\alpha_{/}\right)_{c} :} \& {T^2c} \&\& c \\
		c \&\& Tc \&\& {T^2x} \\
		x \&\& Tx \&\& Tx \\
		\&\& x \&\&\&\& x
		\arrow["{T/f}", from=1-1, to=1-3]
		\arrow["{\alpha_c}", from=1-5, to=1-7]
		\arrow[from=1-5, to=2-5]
		\arrow[from=1-7, to=4-7]
		\arrow[""{name=0, anchor=center, inner sep=0}, from=2-1, to=3-1]
		\arrow[""{name=1, anchor=center, inner sep=0}, from=2-3, to=3-3]
		\arrow["Tf"', from=2-5, to=3-5]
		\arrow["{\alpha_x}", from=2-5, to=4-7]
		\arrow["f", from=3-3, to=4-3]
		\arrow["f"', from=3-5, to=4-7]
		\arrow[between={0.2}{0.8}, maps to, from=0, to=1]
	\end{tikzcd}\cdpunct[8pt]{.}\]
\end{lemma}

\begin{proof}
	The following identity of diagrams implies the result immediately:
	% https://q.uiver.app/#q=WzAsMTUsWzAsMCwiVF4zYyJdLFswLDEsIlReM3giXSxbMCwyLCJUXjJ4Il0sWzAsNCwiVHgiXSxbMyw0LCJ4Il0sWzMsMCwiVGMiXSxbMywzLCJUeCJdLFs0LDEsIlxcZXF1aXYiXSxbNSwwLCJUXjNjIl0sWzUsMSwiVF4zeCJdLFs1LDIsIlReMngiXSxbNSw0LCJUeCJdLFs4LDEsIlR4Il0sWzgsMCwiVGMiXSxbOCw0LCJ4Il0sWzYsNCwiZiJdLFs1LDZdLFszLDQsImYiLDJdLFsyLDMsIlRmIiwyXSxbMCw1LCJUXFxhbHBoYV9jIl0sWzAsMV0sWzEsMiwiVF4yZiIsMl0sWzEsNiwiVFxcYWxwaGFfeCIsMl0sWzEzLDEyXSxbOCwxMywiXFxhbHBoYV97VGN9Il0sWzksMTIsIlxcYWxwaGFfe1R4fSIsMl0sWzgsOV0sWzksMTAsIlReMmYiLDJdLFsxMCwxMSwiVGYiLDJdLFsxMiwxNCwiZiJdLFsxMCwxNCwiXFxhbHBoYV94Il0sWzExLDE0LCJmIiwyXSxbMiw0LCJcXGFscGhhX3giLDJdXQ==
	\[\begin{tikzcd}[ampersand replacement=\&]
		{T^3c} \&\&\& Tc \&\& {T^3c} \&\&\& Tc \\
		{T^3x} \&\&\&\& \equiv \& {T^3x} \&\&\& Tx \\
		{T^2x} \&\&\&\&\& {T^2x} \\
		\&\&\& Tx \\
		Tx \&\&\& x \&\& Tx \&\&\& x
		\arrow["{T\alpha_c}", from=1-1, to=1-4]
		\arrow[from=1-1, to=2-1]
		\arrow[from=1-4, to=4-4]
		\arrow["{\alpha_{Tc}}", from=1-6, to=1-9]
		\arrow[from=1-6, to=2-6]
		\arrow[from=1-9, to=2-9]
		\arrow["{T^2f}"', from=2-1, to=3-1]
		\arrow["{T\alpha_x}"', from=2-1, to=4-4]
		\arrow["{\alpha_{Tx}}"', from=2-6, to=2-9]
		\arrow["{T^2f}"', from=2-6, to=3-6]
		\arrow["f", from=2-9, to=5-9]
		\arrow["Tf"', from=3-1, to=5-1]
		\arrow["{\alpha_x}"', from=3-1, to=5-4]
		\arrow["Tf"', from=3-6, to=5-6]
		\arrow["{\alpha_x}", from=3-6, to=5-9]
		\arrow["f", from=4-4, to=5-4]
		\arrow["f"', from=5-1, to=5-4]
		\arrow["f"', from=5-6, to=5-9]
	\end{tikzcd}\cdpunct[8pt]{.}\]
\end{proof}

As before, there is also a $1$-category $\nInv{0}(C,T,\alpha)$, whose objects are relative involutive objects and whose morphisms are maps $r \colon (x,f) \rightarrow (g,y)$ such that $g \circ Tr = r \circ f$. We can also describe this modified $\nInv{0}$ as a $2$-functor that extends the original functor.

\begin{lemma} \phantom{}
	\begin{enumerate}
		\item There exists a $2$-functor $\nInv{0} \colon \nInv{1}(\nCat{1}) \rightarrow \nCat{1}$ extending the construction of involutive objects along the identity inclusion tritransformation component $i_{\nCat{1}} \colon \nCat{1} \hookrightarrow \nInv{1}(\nCat{1})$.
		\item This further extends to the $2$-functor $\nInv{0} \colon\nInv{1}(\nCat{1}) \rightarrow \nInv{1}(\nCat{1})$ along the forgetful tritransformation component $U_{\nCat{1}} \colon \nInv{1}(\nCat{1}) \rightarrow \nCat{1}$.
		\item The natural $2$-transformation $U$ itself extends to $U \colon \nInv{0} \Rightarrow 1_{\nInv{1}(\nCat{1})}$. However, it doesn't admit any section as before. In particular, this $\nInv{0}$ is not a Monad.
	\end{enumerate}
\end{lemma}

\begin{proof}
	The proof of $(1),(2)$ and $(3)$ mainly follows from a slight modification of the arguments found in \cite[Props.~2.6.2, 2.6.3, and 2.6.5]{yauinvcat2020}.
\end{proof}

We now briefly describe the analogous definitions for inverse structures. Our goal is to provide a better perspective on these definitions that aligns with the microcosm principle and coherence principle.

\begin{definition}[pseudo inverse objects]
	Let $K$ be a bicategory. We define a tuple $(f\colon x \rightarrow y,g\colon y \rightarrow x,\eta,\epsilon)$ to be an inverse object whenever it is an internal adjoint equivalence. That is, the following two triangle identities hold:
	% https://q.uiver.app/#q=WzAsMTgsWzEsMiwieSJdLFswLDAsIngiXSxbMywyLCJ5Il0sWzIsMCwieCJdLFs1LDAsIngiXSxbNiwyLCJ5Il0sWzgsMiwieSJdLFs3LDAsIngiXSxbMCw1LCJ5Il0sWzMsMywieCJdLFsyLDUsInkiXSxbMSwzLCJ4Il0sWzUsNSwieSJdLFs2LDMsIngiXSxbNyw1LCJ5Il0sWzgsMywieCJdLFs0LDQsIlxcZXF1aXYiXSxbNCwxLCJcXGVxdWl2Il0sWzEsMCwiZiIsMl0sWzMsMiwiZiJdLFsxMSwxMCwiZiIsMV0sWzQsNSwiZiIsMl0sWzcsNiwiZiJdLFsxMiwxMywiZyJdLFsxNCwxNSwiZyIsMl0sWzExLDksIlxcdmFyZXBzaWxvbl94Il0sWzEsMywiXFx2YXJlcHNpbG9uX3t4fSJdLFswLDIsIlxcdmFyZXBzaWxvbl97eX0iLDJdLFs1LDYsIlxcdmFyZXBzaWxvbl97eX0iLDJdLFs0LDcsIlxcdmFyZXBzaWxvbl97eH0iXSxbMTIsMTQsIlxcdmFyZXBzaWxvbl95IiwyXSxbMTMsMTUsIlxcdmFyZXBzaWxvbl94Il0sWzAsMywiZyIsMV0sWzgsMTAsIlxcdmFyZXBzaWxvbl95IiwyXSxbOCwxMSwiZyJdLFs0LDYsImYiLDFdLFsxMCw5LCJnIiwyXSxbMTIsMTUsImciLDFdLFsyNiwwLCJcXGV0YSIsMSx7InNob3J0ZW4iOnsic291cmNlIjoyMCwidGFyZ2V0IjoyMH19XSxbMTEsMzMsIlxcZXBzaWxvbiIsMSx7InNob3J0ZW4iOnsic291cmNlIjoyMCwidGFyZ2V0IjoyMH19XSxbMjUsMTAsIlxcZXRhIiwxLHsic2hvcnRlbiI6eyJzb3VyY2UiOjIwLCJ0YXJnZXQiOjIwfX1dLFszLDI3LCJcXGVwc2lsb24iLDEseyJzaG9ydGVuIjp7InNvdXJjZSI6MjAsInRhcmdldCI6MjB9fV1d
	\[\begin{tikzcd}[ampersand replacement=\&]
		x \&\& x \&\&\& x \&\& x \& \\
		\&\&\&\& \equiv \\
		\& y \&\& y \&\&\& y \&\& y \\
		\& x \&\& x \&\&\& x \&\& x \\
		\&\&\&\& \equiv \\
		y \&\& y \&\&\& y \&\& y
		\arrow[""{name=0, anchor=center, inner sep=0}, "{\varepsilon_{x}}", from=1-1, to=1-3]
		\arrow["f"', from=1-1, to=3-2]
		\arrow["f", from=1-3, to=3-4]
		\arrow["{\varepsilon_{x}}", from=1-6, to=1-8]
		\arrow["f"', from=1-6, to=3-7]
		\arrow["f"{description}, from=1-6, to=3-9]
		\arrow["f", from=1-8, to=3-9]
		\arrow["g"{description}, from=3-2, to=1-3]
		\arrow[""{name=1, anchor=center, inner sep=0}, "{\varepsilon_{y}}"', from=3-2, to=3-4]
		\arrow["{\varepsilon_{y}}"', from=3-7, to=3-9]
		\arrow[""{name=2, anchor=center, inner sep=0}, "{\varepsilon_x}", from=4-2, to=4-4]
		\arrow["f"{description}, from=4-2, to=6-3]
		\arrow["{\varepsilon_x}", from=4-7, to=4-9]
		\arrow["g", from=6-1, to=4-2]
		\arrow[""{name=3, anchor=center, inner sep=0}, "{\varepsilon_y}"', from=6-1, to=6-3]
		\arrow["g"', from=6-3, to=4-4]
		\arrow["g", from=6-6, to=4-7]
		\arrow["g"{description}, from=6-6, to=4-9]
		\arrow["{\varepsilon_y}"', from=6-6, to=6-8]
		\arrow["g"', from=6-8, to=4-9]
		\arrow["\eta"{description}, between={0.2}{0.8}, Rightarrow, nfold, from=0, to=3-2]
		\arrow["\epsilon"{description}, between={0.2}{0.8}, Rightarrow, nfold, from=1-3, to=1]
		\arrow["\epsilon"{description}, between={0.2}{0.8}, Rightarrow, nfold, from=4-2, to=3]
		\arrow["\eta"{description}, between={0.2}{0.8}, Rightarrow, nfold, from=2, to=6-3]
	\end{tikzcd}\cdpunct[8pt]{.}\]
\end{definition}

\begin{example}[pseudo inverse objects in $\nCat{1}$]
	A pair of functors $F \colon \cC \rightarrow \cD$ and $G \colon \cD \rightarrow \cC$ specify a weak inverse object whenever they are part of an adjoint equivalence of categories. 
\end{example}

We strongly believe that this shift in perspective of viewing an adjoint equivalence as a coherent weak inverse is important and not merely pedantic. Indeed, assume that data of an adjoint equivalence of categories were genuinely different from the data of an equivalence of categories. In this hypothetical situation, working with equivalence of categories would have been no better than working with monoidal categories without the pentagon or triangle axioms.

\begin{comment}
	\footnote{Infact, this is exactly why the definition of algebraic tricategories in \cite{gurski2006tricat} is better behaved in comparison to original definition of tricategories in \cite{gps1995}. Even if every equivalence can be promoted to an adjoint equivalence, the algebraic tricategories encode this coherence explicitly. In particular, this algebraic definition is the \textit{fully coherent} one in contrast to the original definition.}
\end{comment}

With this idea, we can now work with general adjunctions instead of adjoint equivalences. One could view the data $\eta \colon 1 \Rightarrow GF$ as a laxly categorified version of $f$ being right inverse of $g$ (or $gf=1$), and the similar goes for $\epsilon \colon FG \Rightarrow 1$ as laxly categorified left inverse $f$. Using this viewpoint and \textit{convention}, the data of an adjunction $F \dashv G$ simultaneously expresses that $F$ is a two-sided lax inverse of $G$ and that $G$ is a two-sided colax inverse of $F$.

There is also an analogous notion of $1$-morphisms and $2$-morphisms defining the bicategory of adjunctions in $K$. To avoid notation overload, we will not label the data.

\begin{definition}
	A (lax) $1$-morphism $(x,y,f,g,\eta,\epsilon) \rightarrow (x',y',f',g',\eta',\epsilon')$ consists of
	\begin{enumerate}
		\item two $1$-cells $x \rightarrow x'$ and $y \rightarrow y'$,
		\item two $2$-cells\footnote{These $2$-cells are not mates. One may wish to closely compare this definition with \cite{chengmates2014}. However, the two immediate identities ensures that the $2$nd cell is the inverse of the mate of the first one. In particular, it is a form of Beck-Chevalley condition. Moreover, the colax variant of the $1$-morphisms will yield the dual Beck-Chevalley condition.} as depicted by the following diagram:
		% https://q.uiver.app/#q=WzAsOCxbMCwwLCJ4Il0sWzIsMCwieCciXSxbMCwyLCJ5Il0sWzIsMiwieSciXSxbMywwLCJ5Il0sWzUsMCwieSciXSxbMywyLCJ4Il0sWzUsMiwieCciXSxbMCwyLCJmIiwyXSxbMSwzLCJmJyJdLFswLDFdLFsyLDNdLFsxLDIsIiIsMCx7InNob3J0ZW4iOnsic291cmNlIjoyMCwidGFyZ2V0IjoyMH0sImxldmVsIjoyfV0sWzQsNV0sWzQsNiwiZyIsMl0sWzYsN10sWzUsNywiZyciXSxbNSw2LCIiLDIseyJzaG9ydGVuIjp7InNvdXJjZSI6MjAsInRhcmdldCI6MjB9LCJsZXZlbCI6Mn1dXQ==
		\[\begin{tikzcd}[ampersand replacement=\&]
			x \&\& {x'} \& y \&\& {y'} \\
			\\
			y \&\& {y'} \& x \&\& {x'}
			\arrow[from=1-1, to=1-3]
			\arrow["f"', from=1-1, to=3-1]
			\arrow[between={0.2}{0.8}, Rightarrow, nfold, from=1-3, to=3-1]
			\arrow["{f'}", from=1-3, to=3-3]
			\arrow[from=1-4, to=1-6]
			\arrow["g"', from=1-4, to=3-4]
			\arrow[between={0.2}{0.8}, Rightarrow, nfold, from=1-6, to=3-4]
			\arrow["{g'}", from=1-6, to=3-6]
			\arrow[from=3-1, to=3-3]
			\arrow[from=3-4, to=3-6]
		\end{tikzcd}\cdpunct[8pt]{,}\]
		\item such that the following two identities holds:
		% https://q.uiver.app/#q=WzAsMTIsWzEsMiwieCJdLFszLDIsIngnIl0sWzAsNCwieSJdLFsyLDQsInknIl0sWzAsMCwieSJdLFsyLDAsInknIl0sWzQsMiwiXFxlcXVpdiJdLFs3LDAsInknIl0sWzgsMiwieCciXSxbNyw0LCJ5JyJdLFs1LDAsInkiXSxbNSw0LCJ5Il0sWzAsMiwiZiIsMl0sWzEsMywiZiciXSxbMCwxXSxbMiwzXSxbMSwyLCIiLDAseyJzaG9ydGVuIjp7InNvdXJjZSI6MjAsInRhcmdldCI6MjB9LCJsZXZlbCI6Mn1dLFs0LDVdLFs0LDAsImciLDJdLFs1LDEsImcnIl0sWzUsMCwiIiwyLHsic2hvcnRlbiI6eyJzb3VyY2UiOjIwLCJ0YXJnZXQiOjIwfSwibGV2ZWwiOjJ9XSxbNCwyLCJcXHZhcmVwc2lsb25feSIsMl0sWzcsOCwiZyciXSxbOCw5LCJmJyJdLFs3LDksIlxcdmFyZXBzaWxvbl97eSd9IiwyXSxbMTAsMTEsIlxcdmFyZXBzaWxvbl95IiwyXSxbMTAsN10sWzExLDldLFsxMCw5XSxbMCwyMSwiXFxlcHNpbG9uIiwyLHsic2hvcnRlbiI6eyJzb3VyY2UiOjIwLCJ0YXJnZXQiOjIwfX1dLFs4LDI0LCJcXGVwc2lsb24nIiwyLHsic2hvcnRlbiI6eyJzb3VyY2UiOjIwLCJ0YXJnZXQiOjIwfX1dLFs3LDI4LCIiLDIseyJzaG9ydGVuIjp7InNvdXJjZSI6MjAsInRhcmdldCI6MjB9fV0sWzI4LDExLCIiLDIseyJzaG9ydGVuIjp7InNvdXJjZSI6MjAsInRhcmdldCI6MjB9fV1d
		\[\begin{tikzcd}[ampersand replacement=\&]
			y \&\& {y'} \&\&\& y \&\& {y'} \& \\
			\\
			\& x \&\& {x'} \& \equiv \&\&\&\& {x'} \\
			\\
			y \&\& {y'} \&\&\& y \&\& {y'}
			\arrow[from=1-1, to=1-3]
			\arrow["g"', from=1-1, to=3-2]
			\arrow[""{name=0, anchor=center, inner sep=0}, "{\varepsilon_y}"', from=1-1, to=5-1]
			\arrow[between={0.2}{0.8}, Rightarrow, nfold, from=1-3, to=3-2]
			\arrow["{g'}", from=1-3, to=3-4]
			\arrow[from=1-6, to=1-8]
			\arrow["{\varepsilon_y}"', from=1-6, to=5-6]
			\arrow[""{name=1, anchor=center, inner sep=0}, from=1-6, to=5-8]
			\arrow["{g'}", from=1-8, to=3-9]
			\arrow[""{name=2, anchor=center, inner sep=0}, "{\varepsilon_{y'}}"', from=1-8, to=5-8]
			\arrow[from=3-2, to=3-4]
			\arrow["f"', from=3-2, to=5-1]
			\arrow[between={0.2}{0.8}, Rightarrow, nfold, from=3-4, to=5-1]
			\arrow["{f'}", from=3-4, to=5-3]
			\arrow["{f'}", from=3-9, to=5-8]
			\arrow[from=5-1, to=5-3]
			\arrow[from=5-6, to=5-8]
			\arrow[between={0.2}{0.8}, Rightarrow, nfold, from=1, to=5-6]
			\arrow[between={0.2}{0.8}, Rightarrow, nfold, from=1-8, to=1]
			\arrow["\epsilon"', between={0.2}{0.8}, Rightarrow, nfold, from=3-2, to=0]
			\arrow["{\epsilon'}"', between={0.2}{0.8}, Rightarrow, nfold, from=3-9, to=2]
		\end{tikzcd}\]
		% https://q.uiver.app/#q=WzAsMTIsWzAsMiwieSJdLFsyLDIsInknIl0sWzEsNCwieCJdLFszLDQsIngnIl0sWzEsMCwieCJdLFszLDAsIngnIl0sWzQsMiwiXFxlcXVpdiJdLFs4LDAsIngnIl0sWzgsNCwieCciXSxbNiwwLCJ4Il0sWzYsNCwieCJdLFs1LDIsInkiXSxbMCwyLCJnIiwyXSxbMSwzLCJnJyJdLFswLDFdLFsyLDNdLFsxLDIsIiIsMCx7InNob3J0ZW4iOnsic291cmNlIjoyMCwidGFyZ2V0IjoyMH0sImxldmVsIjoyfV0sWzQsNV0sWzQsMCwiZiIsMl0sWzUsMSwiZiciXSxbNSwwLCIiLDIseyJzaG9ydGVuIjp7InNvdXJjZSI6MjAsInRhcmdldCI6MjB9LCJsZXZlbCI6Mn1dLFs3LDgsIlxcdmFyZXBzaWxvbl97eCd9Il0sWzksN10sWzEwLDhdLFs5LDhdLFs1LDMsIlxcdmFyZXBzaWxvbl97eCd9Il0sWzksMTEsImYiLDJdLFsxMSwxMCwiZyIsMl0sWzksMTAsIlxcdmFyZXBzaWxvbl97eH0iXSxbNywyNCwiIiwyLHsic2hvcnRlbiI6eyJzb3VyY2UiOjIwLCJ0YXJnZXQiOjIwfX1dLFsyNCwxMCwiIiwyLHsic2hvcnRlbiI6eyJzb3VyY2UiOjIwLCJ0YXJnZXQiOjIwfX1dLFsyNSwxLCJcXGV0YSciLDAseyJzaG9ydGVuIjp7InNvdXJjZSI6MjAsInRhcmdldCI6MjB9fV0sWzI4LDExLCJcXGV0YSIsMix7InNob3J0ZW4iOnsic291cmNlIjoyMCwidGFyZ2V0IjoyMH19XV0=
		\[\begin{tikzcd}[ampersand replacement=\&]
			\& x \&\& {x'} \&\&\& x \&\& {x'} \\
			\\
			y \&\& {y'} \&\& \equiv \& y \\
			\\
			\& x \&\& {x'} \&\&\& x \&\& {x'}
			\arrow[from=1-2, to=1-4]
			\arrow["f"', from=1-2, to=3-1]
			\arrow[between={0.2}{0.8}, Rightarrow, nfold, from=1-4, to=3-1]
			\arrow["{f'}", from=1-4, to=3-3]
			\arrow[""{name=0, anchor=center, inner sep=0}, "{\varepsilon_{x'}}", from=1-4, to=5-4]
			\arrow[from=1-7, to=1-9]
			\arrow["f"', from=1-7, to=3-6]
			\arrow[""{name=1, anchor=center, inner sep=0}, "{\varepsilon_{x}}", from=1-7, to=5-7]
			\arrow[""{name=2, anchor=center, inner sep=0}, from=1-7, to=5-9]
			\arrow["{\varepsilon_{x'}}", from=1-9, to=5-9]
			\arrow[from=3-1, to=3-3]
			\arrow["g"', from=3-1, to=5-2]
			\arrow[between={0.2}{0.8}, Rightarrow, nfold, from=3-3, to=5-2]
			\arrow["{g'}", from=3-3, to=5-4]
			\arrow["g"', from=3-6, to=5-7]
			\arrow[from=5-2, to=5-4]
			\arrow[from=5-7, to=5-9]
			\arrow["{\eta'}", between={0.2}{0.8}, Rightarrow, nfold, from=0, to=3-3]
			\arrow["\eta"', between={0.2}{0.8}, Rightarrow, nfold, from=1, to=3-6]
			\arrow[between={0.2}{0.8}, Rightarrow, nfold, from=2, to=5-7]
			\arrow[between={0.2}{0.8}, Rightarrow, nfold, from=1-9, to=2]
		\end{tikzcd}\cdpunct[8pt]{.}\]
	\end{enumerate}
\end{definition}


\begin{definition}
	A $2$-morphism consists of
	\begin{enumerate}
		\item two $2$-cells given as
		% https://q.uiver.app/#q=WzAsNCxbMCwwLCJ4Il0sWzIsMCwieCciXSxbMywwLCJ5Il0sWzUsMCwieSciXSxbMCwxLCIiLDAseyJjdXJ2ZSI6LTJ9XSxbMiwzLCIiLDAseyJjdXJ2ZSI6LTJ9XSxbMCwxLCIiLDIseyJjdXJ2ZSI6Mn1dLFsyLDMsIiIsMix7ImN1cnZlIjoyfV0sWzQsNiwiIiwwLHsic2hvcnRlbiI6eyJzb3VyY2UiOjIwLCJ0YXJnZXQiOjIwfX1dLFs1LDcsIiIsMCx7InNob3J0ZW4iOnsic291cmNlIjoyMCwidGFyZ2V0IjoyMH19XV0=
		\[\begin{tikzcd}[ampersand replacement=\&]
			x \&\& {x'} \& y \&\& {y'}
			\arrow[""{name=0, anchor=center, inner sep=0}, curve={height=-12pt}, from=1-1, to=1-3]
			\arrow[""{name=1, anchor=center, inner sep=0}, curve={height=12pt}, from=1-1, to=1-3]
			\arrow[""{name=2, anchor=center, inner sep=0}, curve={height=-12pt}, from=1-4, to=1-6]
			\arrow[""{name=3, anchor=center, inner sep=0}, curve={height=12pt}, from=1-4, to=1-6]
			\arrow[between={0.2}{0.8}, Rightarrow, nfold, from=0, to=1]
			\arrow[between={0.2}{0.8}, Rightarrow, nfold, from=2, to=3]
		\end{tikzcd}\cdpunct[2pt]{,}\]
		\item such that the following two identities hold:
		% https://q.uiver.app/#q=WzAsMTgsWzAsMCwieCJdLFsyLDAsIngnIl0sWzAsMiwieSJdLFsyLDIsInknIl0sWzAsNCwieSJdLFsyLDQsInknIl0sWzAsNiwieCJdLFsyLDYsIngnIl0sWzMsMSwiXFxlcXVpdiJdLFs0LDAsIngiXSxbNiwwLCJ4JyJdLFs0LDIsInkiXSxbNiwyLCJ5JyJdLFszLDUsIlxcZXF1aXYiXSxbNCw0LCJ5Il0sWzQsNiwieCJdLFs2LDQsInknIl0sWzYsNiwieCciXSxbMCwyLCJmIiwyXSxbMSwzLCJmJyJdLFswLDFdLFsyLDNdLFsxLDIsIiIsMCx7InNob3J0ZW4iOnsic291cmNlIjoyMCwidGFyZ2V0IjoyMH0sImxldmVsIjoyfV0sWzQsNV0sWzQsNiwiZyIsMl0sWzYsN10sWzUsNywiZyciXSxbNSw2LCIiLDIseyJzaG9ydGVuIjp7InNvdXJjZSI6MjAsInRhcmdldCI6MjB9LCJsZXZlbCI6Mn1dLFswLDEsIiIsMCx7ImN1cnZlIjotM31dLFs5LDExLCJmIiwyXSxbMTAsMTIsImYnIl0sWzksMTBdLFsxMSwxMl0sWzExLDEyLCIiLDEseyJjdXJ2ZSI6M31dLFsxMCwxMSwiIiwxLHsic2hvcnRlbiI6eyJzb3VyY2UiOjIwLCJ0YXJnZXQiOjIwfSwibGV2ZWwiOjJ9XSxbNCw1LCIiLDEseyJjdXJ2ZSI6LTN9XSxbMTQsMTUsImciLDJdLFsxNiwxNywiZyciXSxbMTUsMTddLFsxNCwxNl0sWzE1LDE3LCIiLDEseyJjdXJ2ZSI6M31dLFsxNiwxNSwiIiwyLHsic2hvcnRlbiI6eyJzb3VyY2UiOjIwLCJ0YXJnZXQiOjIwfSwibGV2ZWwiOjJ9XSxbMjgsMjAsIiIsMCx7InNob3J0ZW4iOnsic291cmNlIjoyMCwidGFyZ2V0IjoyMH19XSxbMzIsMzMsIiIsMSx7InNob3J0ZW4iOnsic291cmNlIjoyMCwidGFyZ2V0IjoyMH19XSxbMzUsMjMsIiIsMSx7InNob3J0ZW4iOnsic291cmNlIjoyMCwidGFyZ2V0IjoyMH19XSxbMzgsNDAsIiIsMix7InNob3J0ZW4iOnsic291cmNlIjoyMCwidGFyZ2V0IjoyMH19XV0=
		\[\begin{tikzcd}[ampersand replacement=\&]
			x \&\& {x'} \&\& x \&\& {x'} \\
			\&\&\& \equiv \\
			y \&\& {y'} \&\& y \&\& {y'} \\
			\\
			y \&\& {y'} \&\& y \&\& {y'} \\
			\&\&\& \equiv \\
			x \&\& {x'} \&\& x \&\& {x'}
			\arrow[""{name=0, anchor=center, inner sep=0}, from=1-1, to=1-3]
			\arrow[""{name=1, anchor=center, inner sep=0}, curve={height=-18pt}, from=1-1, to=1-3]
			\arrow["f"', from=1-1, to=3-1]
			\arrow[between={0.2}{0.8}, Rightarrow, nfold, from=1-3, to=3-1]
			\arrow["{f'}", from=1-3, to=3-3]
			\arrow[from=1-5, to=1-7]
			\arrow["f"', from=1-5, to=3-5]
			\arrow[between={0.2}{0.8}, Rightarrow, nfold, from=1-7, to=3-5]
			\arrow["{f'}", from=1-7, to=3-7]
			\arrow[from=3-1, to=3-3]
			\arrow[""{name=2, anchor=center, inner sep=0}, from=3-5, to=3-7]
			\arrow[""{name=3, anchor=center, inner sep=0}, curve={height=18pt}, from=3-5, to=3-7]
			\arrow[""{name=4, anchor=center, inner sep=0}, from=5-1, to=5-3]
			\arrow[""{name=5, anchor=center, inner sep=0}, curve={height=-18pt}, from=5-1, to=5-3]
			\arrow["g"', from=5-1, to=7-1]
			\arrow[between={0.2}{0.8}, Rightarrow, nfold, from=5-3, to=7-1]
			\arrow["{g'}", from=5-3, to=7-3]
			\arrow[from=5-5, to=5-7]
			\arrow["g"', from=5-5, to=7-5]
			\arrow[between={0.2}{0.8}, Rightarrow, nfold, from=5-7, to=7-5]
			\arrow["{g'}", from=5-7, to=7-7]
			\arrow[from=7-1, to=7-3]
			\arrow[""{name=6, anchor=center, inner sep=0}, from=7-5, to=7-7]
			\arrow[""{name=7, anchor=center, inner sep=0}, curve={height=18pt}, from=7-5, to=7-7]
			\arrow[between={0.2}{0.8}, Rightarrow, nfold, from=1, to=0]
			\arrow[between={0.2}{0.8}, Rightarrow, nfold, from=2, to=3]
			\arrow[between={0.2}{0.8}, Rightarrow, nfold, from=5, to=4]
			\arrow[between={0.2}{0.8}, Rightarrow, nfold, from=6, to=7]
		\end{tikzcd}\cdpunct[8pt]{.}\]
	\end{enumerate}
\end{definition}

One objection the reader may have is that we never discussed the composites of adjoints in the context of our proposed \emph{bicategory of adjunctions}. To resolve this discrepancy, we shall illustrate the idea one dimension lower with a straightforward lemma.

\begin{lemma}
	There exists a double category of inverses in $\cC$ defined by the following structure:
	\begin{description}
		\item[Objects - ] are precisely the objects of $\cC$ itself,
		\item[Horizontal $1$-cells - ] are precisely the morphisms of $\cC$,
		\item[Vertical $1$-cells - ] are precisely the invertible arrows in $\cC$,
		\item[Square $2$-cells  - ] are commuting squares.
	\end{description}
\end{lemma}

From this double category, one can extract the usual $1$-category of inverses, whose objects are invertible pairs of morphisms $(f,f^{-1})$, and whose morphisms are commuting square.

This suggests that our bicategory of adjunction in $K$ should itself be a part of some higher double categorical structure, such as a Verity double bicategory \cite{morton2009doublebicat} or even an intercategory \cite{grandis2017intercatframework}. However, developing this structure is outside the scope of this paper, and we shall not pursue it further here.

Analogous to the involution case, there is also a \emph{$1$-category of relative inverse objects}. Furthermore, adjunctions admit slicing using these inverse objects.

\begin{definition}[inverse objects]
	Let $(F\colon \cC \rightarrow \cD,G\colon \cD \rightarrow \cC,\eta,\epsilon)$ be an adjunction. We define a relative inverse object to be the pair $(f \colon Fc \rightarrow d , g \colon Gd \rightarrow c)$, such that the morphisms satisfy the following identities:
	% https://q.uiver.app/#q=WzAsOCxbMCwwLCJHKEZjKSJdLFsyLDAsIkcoZCkiXSxbMCwxLCJjIl0sWzMsMCwiRihHZCkiXSxbNSwwLCJGKGMpIl0sWzMsMSwiZCJdLFsyLDEsImMiXSxbNSwxLCJkIl0sWzAsMSwiRyhmKSJdLFsyLDAsIlxcZXRhX2MiXSxbMyw0LCJGKGcpIl0sWzMsNSwiXFxlcHNpbG9uX2QiLDJdLFsxLDYsImciXSxbMiw2LCIxX2MiLDJdLFs0LDcsImYiXSxbNSw3LCIxX2QiLDJdXQ==
	\[\begin{tikzcd}[ampersand replacement=\&]
		{G(Fc)} \&\& {G(d)} \& {F(Gd)} \&\& {F(c)} \\
		c \&\& c \& d \&\& d
		\arrow["{G(f)}", from=1-1, to=1-3]
		\arrow["g", from=1-3, to=2-3]
		\arrow["{F(g)}", from=1-4, to=1-6]
		\arrow["{\epsilon_d}"', from=1-4, to=2-4]
		\arrow["f", from=1-6, to=2-6]
		\arrow["{\eta_c}", from=2-1, to=1-1]
		\arrow["{1_c}"', from=2-1, to=2-3]
		\arrow["{1_d}"', from=2-4, to=2-6]
	\end{tikzcd}\cdpunct[8pt]{.}\]
	
\end{definition}

\begin{lemma}
	Let $F \dashv G$ be an adjunction (not necessarily equivalence) between $\cC$ and $\cD$. And let $(f \colon Fc \rightarrow d, g \colon Gd \rightarrow c)$ be the data of an inverse object relative to this adjunction. Then, there exists a canonical sliced adjunction:
	% https://q.uiver.app/#q=WzAsMTIsWzEsMiwiXFxjQy9jIl0sWzMsMiwiXFxjRC9kIl0sWzEsMSwiYyJdLFsxLDAsImNfMSJdLFszLDEsIkYoYykiXSxbMywwLCJGKGNfMSkiXSxbNCwxLCJkIl0sWzMsMywiZF8xIl0sWzMsNCwiZCJdLFsxLDMsIkcoZF8xKSJdLFsxLDQsIkcoZCkiXSxbMCw0LCJjIl0sWzAsMSwiRl97L2Z9IiwwLHsib2Zmc2V0IjotMn1dLFsxLDAsIkdfey9nfSIsMCx7Im9mZnNldCI6LTJ9XSxbMywyXSxbNCw2LCJmIl0sWzUsNF0sWzEwLDExLCJnIl0sWzksMTBdLFs3LDhdLFsxMiwxMywiIiwwLHsibGV2ZWwiOjEsInN0eWxlIjp7Im5hbWUiOiJhZGp1bmN0aW9uIn19XSxbMTksMTgsIiIsMCx7InNob3J0ZW4iOnsic291cmNlIjoyMCwidGFyZ2V0IjoyMH0sImxldmVsIjoxLCJzdHlsZSI6eyJ0YWlsIjp7Im5hbWUiOiJtYXBzIHRvIn19fV0sWzE0LDE2LCIiLDAseyJzaG9ydGVuIjp7InNvdXJjZSI6MjAsInRhcmdldCI6MjB9LCJsZXZlbCI6MSwic3R5bGUiOnsidGFpbCI6eyJuYW1lIjoibWFwcyB0byJ9fX1dXQ==
	\[\begin{tikzcd}[ampersand replacement=\&]
		\& {c_1} \&\& {F(c_1)} \& \\
		\& c \&\& {F(c)} \& d \\
		\& {\cC/c} \&\& {\cD/d} \\
		\& {G(d_1)} \&\& {d_1} \\
		c \& {G(d)} \&\& d
		\arrow[""{name=0, anchor=center, inner sep=0}, from=1-2, to=2-2]
		\arrow[""{name=1, anchor=center, inner sep=0}, from=1-4, to=2-4]
		\arrow["f", from=2-4, to=2-5]
		\arrow[""{name=2, anchor=center, inner sep=0}, "{F_{/f}}", shift left=2, from=3-2, to=3-4]
		\arrow[""{name=3, anchor=center, inner sep=0}, "{G_{/g}}", shift left=2, from=3-4, to=3-2]
		\arrow[""{name=4, anchor=center, inner sep=0}, from=4-2, to=5-2]
		\arrow[""{name=5, anchor=center, inner sep=0}, from=4-4, to=5-4]
		\arrow["g", from=5-2, to=5-1]
		\arrow[between={0.2}{0.8}, maps to, from=0, to=1]
		\arrow["\dashv"{anchor=center, rotate=-90}, draw=none, from=2, to=3]
		\arrow[between={0.2}{0.8}, maps to, from=5, to=4]
	\end{tikzcd}\cdpunct[8pt]{.}\]
	Here, the action on morphisms is the one induced by $F$ and $G$.
\end{lemma}

\begin{proof}
	Let us denote the new adjoint pair using the tuple $(F_{/f}, G_{/g}, \eta_{/}, \epsilon_{/})$. Define $\eta_{/}$ and $\epsilon_{/}$ using the following two diagrams:
	% https://q.uiver.app/#q=WzAsMTQsWzEsMCwiY18xIl0sWzIsMywiYyJdLFszLDAsIkdGKGNfMSkiXSxbMywxLCJHRihjKSJdLFszLDIsIkdkIl0sWzEsMiwiYyJdLFswLDAsIntcXGxlZnQoXFxldGFfey99XFxyaWdodCl9X3tcXGFscGhhfSA6Il0sWzQsMCwie1xcbGVmdChcXGVwc2lsb25fey99XFxyaWdodCl9X3tcXGJldGF9IDoiXSxbNSwwLCJGRyhkXzEpIl0sWzUsMSwiRkcoZCkiXSxbNSwyLCJGYyJdLFs2LDMsImQiXSxbNywyLCJkIl0sWzcsMCwiZF8xIl0sWzQsMSwiZyJdLFszLDQsIkdmIl0sWzIsMywiR0YoXFxhbHBoYSkiXSxbMCwyLCJcXGV0YV97Y18xfSJdLFswLDUsIlxcYWxwaGEiLDJdLFs1LDMsIlxcZXRhX2MiXSxbNSwxLCIxX3tjfSIsMl0sWzgsOSwiRkcoXFxiZXRhKSIsMl0sWzEzLDEyLCJcXGJldGEiXSxbOCwxMywiXFxlcHNpbG9uX3tkXzF9Il0sWzksMTAsIkYoZykiLDJdLFsxMCwxMSwiZiIsMl0sWzEyLDExLCIxX3tkfSJdLFs5LDEyLCJcXGVwc2lsb25fe2R9Il1d
	\[\begin{tikzcd}[ampersand replacement=\&]
		{{\left(\eta_{/}\right)}_{\alpha} :} \& {c_1} \&\& {GF(c_1)} \& {{\left(\epsilon_{/}\right)}_{\beta} :} \& {FG(d_1)} \&\& {d_1} \\
		\&\&\& {GF(c)} \&\& {FG(d)} \\
		\& c \&\& Gd \&\& Fc \&\& d \\
		\&\& c \&\&\&\& d
		\arrow["{\eta_{c_1}}", from=1-2, to=1-4]
		\arrow["\alpha"', from=1-2, to=3-2]
		\arrow["{GF(\alpha)}", from=1-4, to=2-4]
		\arrow["{\epsilon_{d_1}}", from=1-6, to=1-8]
		\arrow["{FG(\beta)}"', from=1-6, to=2-6]
		\arrow["\beta", from=1-8, to=3-8]
		\arrow["Gf", from=2-4, to=3-4]
		\arrow["{F(g)}"', from=2-6, to=3-6]
		\arrow["{\epsilon_{d}}", from=2-6, to=3-8]
		\arrow["{\eta_c}", from=3-2, to=2-4]
		\arrow["{1_{c}}"', from=3-2, to=4-3]
		\arrow["g", from=3-4, to=4-3]
		\arrow["f"', from=3-6, to=4-7]
		\arrow["{1_{d}}", from=3-8, to=4-7]
	\end{tikzcd}\cdpunct[8pt]{.}\]
	The relevant triangle identities are then immediate by pasting these two diagrams appropriately. For instance, one of the triangle identity is given by the following commuting diagram:
	% https://q.uiver.app/#q=WzAsMTAsWzAsMCwiRihjXzEpIl0sWzEsMywiRihjKSJdLFsyLDAsIkZHRihjXzEpIl0sWzIsMSwiRkdGKGMpIl0sWzIsMiwiRkdkIl0sWzAsMiwiRihjKSJdLFsyLDQsImQiXSxbNCwzLCJkIl0sWzQsMiwiRihjKSJdLFs0LDAsIkYoY18xKSJdLFs0LDEsIkYoZykiXSxbMyw0LCJGRyhmKSJdLFsyLDMsIkZHRihcXGFscGhhKSJdLFswLDIsIkZcXGV0YV97Y18xfSJdLFswLDUsIkYoXFxhbHBoYSkiLDJdLFs1LDMsIkZcXGV0YV9jIl0sWzUsMSwiMV97RihjKX0iLDJdLFsxLDYsImYiLDJdLFs0LDcsIlxcZXBzaWxvbl97ZH0iXSxbMyw4LCJcXGVwc2lsb25fe0ZjfSJdLFsyLDksIlxcZXBzaWxvbl97RmNfMX0iXSxbOSw4LCJGKFxcYWxwaGEpIl0sWzgsNywiZiJdLFs3LDYsIjFfe2R9Il1d
	\[\begin{tikzcd}[ampersand replacement=\&]
		{F(c_1)} \&\& {FGF(c_1)} \&\& {F(c_1)} \\
		\&\& {FGF(c)} \\
		{F(c)} \&\& FGd \&\& {F(c)} \\
		\& {F(c)} \&\&\& d \\
		\&\& d
		\arrow["{F\eta_{c_1}}", from=1-1, to=1-3]
		\arrow["{F(\alpha)}"', from=1-1, to=3-1]
		\arrow["{\epsilon_{Fc_1}}", from=1-3, to=1-5]
		\arrow["{FGF(\alpha)}", from=1-3, to=2-3]
		\arrow["{F(\alpha)}", from=1-5, to=3-5]
		\arrow["{FG(f)}", from=2-3, to=3-3]
		\arrow["{\epsilon_{Fc}}", from=2-3, to=3-5]
		\arrow["{F\eta_c}", from=3-1, to=2-3]
		\arrow["{1_{F(c)}}"', from=3-1, to=4-2]
		\arrow["{F(g)}", from=3-3, to=4-2]
		\arrow["{\epsilon_{d}}", from=3-3, to=4-5]
		\arrow["f", from=3-5, to=4-5]
		\arrow["f"', from=4-2, to=5-3]
		\arrow["{1_{d}}", from=4-5, to=5-3]
	\end{tikzcd}\cdpunct[8pt]{.}\]
\end{proof}

At this stage, we have provided ample motivation and theory to view dualities and inverses through the lens of the microcosm principle. However, the original statement holds between two arbitrary dimensions $n$ and $n+1$. Therefore, a similar theory of dualities and inverses should also hold in one dimension higher. In particular, we want to define relative pseudoinvolutions and pseudoinverses.

\begin{comment}
	To this end, we first note the following counterpart of lemma (see):
	
	\begin{lemma}
		Let $T$ be a tricategory. Let $H : K \rightarrow K$ be a $1$-cell in $T$ equipped with an adjoint equivalence $\alpha : H \otimes H \iff 1_{K} : \alpha^{.}$. The following modifications are derivable from each other through mates (and some coherence constraints):
		\begin{enumerate}
			\item % https://q.uiver.app/#q=WzAsMTMsWzIsMiwiSyJdLFszLDAsIksiXSxbMyw0LCJLIl0sWzAsMiwiSyJdLFs1LDIsIksiXSxbOCwwLCJLIl0sWzgsNCwiSyJdLFs0LDMsIlxcUnJpZ2h0YXJyb3dee1xccGhpfSJdLFsyLDcsIkggXFxvdGltZXMgMV9LIl0sWzIsNSwiSCBcXG90aW1lcyAoSCBcXG90aW1lcyBIKSAiXSxbNiw1LCIoSCBcXG90aW1lcyBIKSBcXG90aW1lcyBIIl0sWzYsNywiMV9LIFxcb3RpbWVzIEgiXSxbNCw3LCJIIl0sWzEsMiwiSCJdLFswLDIsIjFfSyIsMl0sWzMsMSwiMV9LIl0sWzQsNiwiSCIsMV0sWzQsNSwiMV9LIl0sWzUsNiwiSCJdLFswLDEsIkgiLDFdLFszLDAsIkgiLDJdLFs4LDksIjEgXFxvdGltZXMgXFxhbHBoYV57Ln0iXSxbOSwxMCwiYV57Ln0iXSxbMTAsMTEsIlxcYWxwaGEgXFxvdGltZXMgMSJdLFs4LDEyLCJyIiwyXSxbMTEsMTIsImwiXSxbMywyLCJIIiwyXSxbMSwxNCwiXFxhbHBoYSIsMCx7InNob3J0ZW4iOnsic291cmNlIjoyMCwidGFyZ2V0IjoyMH19XSxbMTUsMCwiXFxhbHBoYV57Ln0iLDIseyJzaG9ydGVuIjp7InNvdXJjZSI6MjAsInRhcmdldCI6MjB9fV0sWzUsMTYsInIiLDAseyJzaG9ydGVuIjp7InNvdXJjZSI6MjAsInRhcmdldCI6MzB9fV0sWzIyLDEyLCJcXHBoaSIsMCx7InNob3J0ZW4iOnsic291cmNlIjoyMCwidGFyZ2V0IjoyMH19XSxbMCwyNiwibCIsMix7InNob3J0ZW4iOnsic291cmNlIjoyMCwidGFyZ2V0IjoyMH19XV0=
			\[\begin{tikzcd}[ampersand replacement=\&]
				\&\&\& K \&\&\&\&\& K \\
				\\
				K \&\& K \&\&\& K \\
				\&\&\&\& {\Rrightarrow^{\phi}} \\
				\&\&\& K \&\&\&\&\& K \\
				\&\& {H \otimes (H \otimes H) } \&\&\&\& {(H \otimes H) \otimes H} \\
				\\
				\&\& {H \otimes 1_K} \&\& H \&\& {1_K \otimes H}
				\arrow["H", from=1-4, to=5-4]
				\arrow["H", from=1-9, to=5-9]
				\arrow[""{name=0, anchor=center, inner sep=0}, "{1_K}", from=3-1, to=1-4]
				\arrow["H"', from=3-1, to=3-3]
				\arrow[""{name=1, anchor=center, inner sep=0}, "H"', from=3-1, to=5-4]
				\arrow["H"{description}, from=3-3, to=1-4]
				\arrow[""{name=2, anchor=center, inner sep=0}, "{1_K}"', from=3-3, to=5-4]
				\arrow["{1_K}", from=3-6, to=1-9]
				\arrow[""{name=3, anchor=center, inner sep=0}, "H"{description}, from=3-6, to=5-9]
				\arrow[""{name=4, anchor=center, inner sep=0}, "{a^{.}}", from=6-3, to=6-7]
				\arrow["{\alpha \otimes 1}", from=6-7, to=8-7]
				\arrow["{1 \otimes \alpha^{.}}", from=8-3, to=6-3]
				\arrow["r"', from=8-3, to=8-5]
				\arrow["l", from=8-7, to=8-5]
				\arrow["\alpha", between={0.2}{0.8}, Rightarrow, nfold, from=1-4, to=2]
				\arrow["r", between={0.2}{0.7}, Rightarrow, nfold, from=1-9, to=3]
				\arrow["{\alpha^{.}}"', between={0.2}{0.8}, Rightarrow, nfold, from=0, to=3-3]
				\arrow["l"', between={0.2}{0.8}, Rightarrow, nfold, from=3-3, to=1]
				\arrow["\phi", between={0.2}{0.8}, Rightarrow, nfold, from=4, to=8-5]
			\end{tikzcd}\]
			\item % https://q.uiver.app/#q=WzAsMTQsWzEsMiwiSyJdLFszLDAsIksiXSxbMyw0LCJLIl0sWzAsNCwiSyJdLFs1LDQsIksiXSxbNywyLCJLIl0sWzgsNCwiSyJdLFs0LDMsIlxcUnJpZ2h0YXJyb3dee1xcemV0YX0iXSxbMiw3LCJIIFxcb3RpbWVzIDFfSyJdLFsyLDUsIkggXFxvdGltZXMgKEggXFxvdGltZXMgSCkgIl0sWzYsNSwiKEggXFxvdGltZXMgSCkgXFxvdGltZXMgSCJdLFs2LDcsIjFfSyBcXG90aW1lcyBIIl0sWzQsOCwiSCJdLFs1LDAsIksiXSxbMSwyLCJIIl0sWzAsMiwiMV9LIiwyXSxbNCw2LCJIIiwxXSxbNCw1LCIxX0siXSxbNSw2LCJIIl0sWzAsMSwiSCIsMV0sWzMsMCwiSCIsMl0sWzksOCwiMSBcXG90aW1lcyBcXGFscGhhIiwyXSxbOSwxMCwiYV57Ln0iXSxbMTAsMTEsIlxcYWxwaGEgXFxvdGltZXMgMSJdLFs4LDEyLCJyIiwyXSxbMTEsMTIsImwiXSxbMywyLCJIIiwyXSxbNCwxM10sWzEzLDVdLFsxMCw4LCJcXHpldGEiLDAseyJzaG9ydGVuIjp7InNvdXJjZSI6MjAsInRhcmdldCI6MjB9LCJsZXZlbCI6Mn1dLFsxLDE1LCJcXGFscGhhIiwwLHsic2hvcnRlbiI6eyJzb3VyY2UiOjIwLCJ0YXJnZXQiOjIwfX1dLFs1LDE2LCJyIiwwLHsic2hvcnRlbiI6eyJzb3VyY2UiOjIwLCJ0YXJnZXQiOjMwfX1dLFswLDI2LCJsIiwyLHsic2hvcnRlbiI6eyJzb3VyY2UiOjIwLCJ0YXJnZXQiOjIwfX1dLFsxMywxNywiXFxhbHBoYSIsMCx7InNob3J0ZW4iOnsic291cmNlIjoyMCwidGFyZ2V0IjoyMH19XV0=
			\[\begin{tikzcd}[ampersand replacement=\&]
				\&\&\& K \&\& K \&\&\& \\
				\\
				\& K \&\&\&\&\&\& K \\
				\&\&\&\& {\Rrightarrow^{\zeta}} \\
				K \&\&\& K \&\& K \&\&\& K \\
				\&\& {H \otimes (H \otimes H) } \&\&\&\& {(H \otimes H) \otimes H} \\
				\\
				\&\& {H \otimes 1_K} \&\&\&\& {1_K \otimes H} \\
				\&\&\&\& H
				\arrow["H", from=1-4, to=5-4]
				\arrow[from=1-6, to=3-8]
				\arrow["H"{description}, from=3-2, to=1-4]
				\arrow[""{name=0, anchor=center, inner sep=0}, "{1_K}"', from=3-2, to=5-4]
				\arrow["H", from=3-8, to=5-9]
				\arrow["H"', from=5-1, to=3-2]
				\arrow[""{name=1, anchor=center, inner sep=0}, "H"', from=5-1, to=5-4]
				\arrow[from=5-6, to=1-6]
				\arrow[""{name=2, anchor=center, inner sep=0}, "{1_K}", from=5-6, to=3-8]
				\arrow[""{name=3, anchor=center, inner sep=0}, "H"{description}, from=5-6, to=5-9]
				\arrow["{a^{.}}", from=6-3, to=6-7]
				\arrow["{1 \otimes \alpha}"', from=6-3, to=8-3]
				\arrow["\zeta", between={0.2}{0.8}, Rightarrow, nfold, from=6-7, to=8-3]
				\arrow["{\alpha \otimes 1}", from=6-7, to=8-7]
				\arrow["r"', from=8-3, to=9-5]
				\arrow["l", from=8-7, to=9-5]
				\arrow["\alpha", between={0.2}{0.8}, Rightarrow, nfold, from=1-4, to=0]
				\arrow["\alpha", between={0.2}{0.8}, Rightarrow, nfold, from=1-6, to=2]
				\arrow["l"', between={0.2}{0.8}, Rightarrow, nfold, from=3-2, to=1]
				\arrow["r", between={0.2}{0.7}, Rightarrow, nfold, from=3-8, to=3]
			\end{tikzcd}\]
			\item % https://q.uiver.app/#q=WzAsMTMsWzAsMCwiSyJdLFszLDIsIksiXSxbMCw0LCJLIl0sWzUsNCwiSyJdLFs2LDIsIksiXSxbOCwyLCJLIl0sWzQsMiwiXFxScmlnaHRhcnJvd157XFxiYXJ7XFxwaGl9fSJdLFsyLDcsIkggXFxvdGltZXMgMV9LIl0sWzIsNSwiSCBcXG90aW1lcyAoSCBcXG90aW1lcyBIKSAiXSxbNiw1LCIoSCBcXG90aW1lcyBIKSBcXG90aW1lcyBIIl0sWzYsNywiMV9LIFxcb3RpbWVzIEgiXSxbNCw4LCJIIl0sWzUsMCwiSyJdLFswLDEsIjFfSyIsMl0sWzMsNSwiSCIsMV0sWzMsNCwiMV9LIiwxXSxbNCw1LCJIIiwxXSxbMiwwLCJIIiwyXSxbOCw3LCIxIFxcb3RpbWVzIFxcYWxwaGEiLDJdLFs5LDgsImEiLDJdLFsxMCw5LCJcXGFscGhhXnsufSBcXG90aW1lcyAxIiwyXSxbNywxMSwiciIsMl0sWzEwLDExLCJsIl0sWzIsMSwiSCIsMl0sWzMsMTIsIkgiXSxbMTIsNCwiSCIsMV0sWzEyLDUsIjFfSyJdLFsxMCw4LCJcXGJhcntcXHBoaX0iLDIseyJzaG9ydGVuIjp7InNvdXJjZSI6MjAsInRhcmdldCI6MjB9LCJsZXZlbCI6Mn1dLFs0LDE0LCJyIiwwLHsic2hvcnRlbiI6eyJzb3VyY2UiOjIwLCJ0YXJnZXQiOjMwfX1dLFswLDIzLCJsIiwyLHsic2hvcnRlbiI6eyJzb3VyY2UiOjIwLCJ0YXJnZXQiOjIwfX1dLFsyNiw0LCJcXGFscGhhXnsufSIsMCx7InNob3J0ZW4iOnsic291cmNlIjoyMH19XSxbMTIsMTUsIlxcYWxwaGEiLDAseyJzaG9ydGVuIjp7InNvdXJjZSI6MjAsInRhcmdldCI6MjB9fV1d
			\[\begin{tikzcd}[ampersand replacement=\&]
				K \&\&\&\&\& K \&\&\& \\
				\\
				\&\&\& K \& {\Rrightarrow^{\bar{\phi}}} \&\& K \&\& K \\
				\\
				K \&\&\&\&\& K \\
				\&\& {H \otimes (H \otimes H) } \&\&\&\& {(H \otimes H) \otimes H} \\
				\\
				\&\& {H \otimes 1_K} \&\&\&\& {1_K \otimes H} \\
				\&\&\&\& H
				\arrow["{1_K}"', from=1-1, to=3-4]
				\arrow["H"{description}, from=1-6, to=3-7]
				\arrow[""{name=0, anchor=center, inner sep=0}, "{1_K}", from=1-6, to=3-9]
				\arrow["H"{description}, from=3-7, to=3-9]
				\arrow["H"', from=5-1, to=1-1]
				\arrow[""{name=1, anchor=center, inner sep=0}, "H"', from=5-1, to=3-4]
				\arrow["H", from=5-6, to=1-6]
				\arrow[""{name=2, anchor=center, inner sep=0}, "{1_K}"{description}, from=5-6, to=3-7]
				\arrow[""{name=3, anchor=center, inner sep=0}, "H"{description}, from=5-6, to=3-9]
				\arrow["{1 \otimes \alpha}"', from=6-3, to=8-3]
				\arrow["a"', from=6-7, to=6-3]
				\arrow["r"', from=8-3, to=9-5]
				\arrow["{\bar{\phi}}"', between={0.2}{0.8}, Rightarrow, nfold, from=8-7, to=6-3]
				\arrow["{\alpha^{.} \otimes 1}"', from=8-7, to=6-7]
				\arrow["l", from=8-7, to=9-5]
				\arrow["l"', between={0.2}{0.8}, Rightarrow, nfold, from=1-1, to=1]
				\arrow["{\alpha^{.}}", between={0.2}{1}, Rightarrow, nfold, from=0, to=3-7]
				\arrow["\alpha", between={0.2}{0.8}, Rightarrow, nfold, from=1-6, to=2]
				\arrow["r", between={0.2}{0.7}, Rightarrow, nfold, from=3-7, to=3]
			\end{tikzcd}\]
			\item 
		\end{enumerate}
	\end{lemma}
\end{comment}


\subsection{\tpdf{$2$-dimensional variants}} \label{sub:2diminv}

To obtain the definition of $2$-dimensional variant of involutive objects, it is best to utilize a categorification of the lax compatible form of $1$-dimensional involutive objects. A somewhat similar definition can be found in \cite[Def.~2.1]{shulman2018contenrich} for the special case of $T = \BiCat$.

\begin{definition}
	A weak $2$-involution in a tricategory $T$ consists of
	\begin{enumerate}
		\item a $1$-cell $H \colon K \rightarrow K$,
		\item an adjoint equivalence $\alpha \colon H \otimes H \iff 1_K \colon \alpha^{.} $,
		\item a $3$-cell $\zeta$ as depicted by the following diagram:
		% https://q.uiver.app/#q=WzAsMTQsWzAsMCwiSyJdLFsyLDAsIksiXSxbMiwyLCJLIl0sWzAsMiwiSyJdLFs0LDIsIksiXSxbNiwwLCJLIl0sWzYsMiwiSyJdLFszLDEsIlxcUnJpZ2h0YXJyb3dee1xcemV0YX0iXSxbNSw1LCJIIFxcb3RpbWVzIDFfSyJdLFs1LDMsIkggXFxvdGltZXMgKEggXFxvdGltZXMgSCkgIl0sWzEsMywiKEggXFxvdGltZXMgSCkgXFxvdGltZXMgSCJdLFsxLDUsIjFfSyBcXG90aW1lcyBIIl0sWzMsNiwiSCJdLFs0LDAsIksiXSxbMSwyLCJIIl0sWzAsMiwiMV9LIiwxXSxbNCw2LCJIIiwxXSxbNCw1LCIxX0siLDFdLFs1LDYsIkgiXSxbMCwxLCJIIl0sWzMsMCwiSCJdLFs5LDgsIjEgXFxvdGltZXMgXFxhbHBoYSJdLFsxMCw5LCJhIl0sWzEwLDExLCJcXGFscGhhIFxcb3RpbWVzIDEiLDJdLFs4LDEyLCJyIl0sWzExLDEyLCJsIiwyXSxbMywyLCJIIiwyXSxbNCwxMywiSCJdLFsxMyw1LCJIIl0sWzExLDksIlxcemV0YSIsMCx7InNob3J0ZW4iOnsic291cmNlIjoyMCwidGFyZ2V0IjoyMH0sImxldmVsIjoyfV0sWzEsMTUsIlxcYWxwaGEiLDAseyJzaG9ydGVuIjp7InNvdXJjZSI6MjAsInRhcmdldCI6MjB9fV0sWzUsMTYsInIiLDAseyJzaG9ydGVuIjp7InNvdXJjZSI6MjAsInRhcmdldCI6MzB9fV0sWzAsMjYsImwiLDIseyJzaG9ydGVuIjp7InNvdXJjZSI6MjAsInRhcmdldCI6MjB9fV0sWzEzLDE3LCJcXGFscGhhIiwwLHsic2hvcnRlbiI6eyJzb3VyY2UiOjIwLCJ0YXJnZXQiOjIwfX1dXQ==
		\[\begin{tikzcd}[ampersand replacement=\&,cramped]
			K \&\& K \&\& K \&\& K \\
			\&\&\& {\Rrightarrow^{\zeta}} \\
			K \&\& K \&\& K \&\& K \\
			\& {(H \otimes H) \otimes H} \&\&\&\& {H \otimes (H \otimes H) } \\
			\\
			\& {1_K \otimes H} \&\&\&\& {H \otimes 1_K} \\
			\&\&\& H
			\arrow["H", from=1-1, to=1-3]
			\arrow[""{name=0, anchor=center, inner sep=0}, "{1_K}"{description}, from=1-1, to=3-3]
			\arrow["H", from=1-3, to=3-3]
			\arrow["H", from=1-5, to=1-7]
			\arrow["H", from=1-7, to=3-7]
			\arrow["H", from=3-1, to=1-1]
			\arrow[""{name=1, anchor=center, inner sep=0}, "H"', from=3-1, to=3-3]
			\arrow["H", from=3-5, to=1-5]
			\arrow[""{name=2, anchor=center, inner sep=0}, "{1_K}"{description}, from=3-5, to=1-7]
			\arrow[""{name=3, anchor=center, inner sep=0}, "H"{description}, from=3-5, to=3-7]
			\arrow["a", from=4-2, to=4-6]
			\arrow["{\alpha \otimes 1}"', from=4-2, to=6-2]
			\arrow["{1 \otimes \alpha}", from=4-6, to=6-6]
			\arrow["\zeta", between={0.2}{0.8}, Rightarrow, nfold, from=6-2, to=4-6]
			\arrow["l"', from=6-2, to=7-4]
			\arrow["r", from=6-6, to=7-4]
			\arrow["l"', between={0.2}{0.8}, Rightarrow, nfold, from=1-1, to=1]
			\arrow["\alpha", between={0.2}{0.8}, Rightarrow, nfold, from=1-3, to=0]
			\arrow["\alpha", between={0.2}{0.8}, Rightarrow, nfold, from=1-5, to=2]
			\arrow["r", between={0.2}{0.7}, Rightarrow, nfold, from=1-7, to=3]
		\end{tikzcd}\cdpunct[8pt]{,}\]
		\item the $3$-cell $\zeta$ satisfy the weak $2$-involution axiom (\cref{app:weak2invdef}).
		
	\end{enumerate}
\end{definition}

Although the morphisms of weak $2$-involutions are not directly relevant for our purposes, we record their definitions in \cref{appsub:weak2inv}. We now define the relative pseudo involutive objects as suggested by the microcosm principle.

\begin{definition}
	Let $(K,H,\alpha,\zeta)$ be a weak $2$-involution in the tricategory $\BiCat$. We define a relative lax involutive object by
	\begin{enumerate}
		\item a $1$-cell $f \colon Hx \rightarrow x \in B$,
		\item a $2$-cell depicted by the following diagram:
		% https://q.uiver.app/#q=WzAsNCxbMCwwLCJIXjJ4Il0sWzIsMCwiSHgiXSxbMiwyLCJ4Il0sWzAsMiwieCJdLFsxLDIsImYiXSxbMCwxLCJIZiJdLFswLDMsIlxcYWxwaGFfeCIsMl0sWzMsMiwiXFx2YXJlcHNpbG9uX3giLDJdLFsxLDMsIlxccGhpIiwwLHsic2hvcnRlbiI6eyJzb3VyY2UiOjIwLCJ0YXJnZXQiOjIwfSwibGV2ZWwiOjJ9XV0=
		\[\begin{tikzcd}[ampersand replacement=\&,cramped]
			{H^2x} \&\& Hx \\
			\\
			x \&\& x
			\arrow["Hf", from=1-1, to=1-3]
			\arrow["{\alpha_x}"', from=1-1, to=3-1]
			\arrow["\phi", between={0.2}{0.8}, Rightarrow, nfold, from=1-3, to=3-1]
			\arrow["f", from=1-3, to=3-3]
			\arrow["{\varepsilon_x}"', from=3-1, to=3-3]
		\end{tikzcd}\cdpunct[8pt]{,}\]
		\item such that the $2$-cell satisfy the following identity:
		% https://q.uiver.app/#q=WzAsMTIsWzYsMCwiSF4zeCJdLFs4LDAsIkheMngiXSxbOCwyLCJIeCJdLFsxMCwyLCJ4Il0sWzIsMCwiSF4yeCJdLFs0LDAsIkh4Il0sWzQsMiwieCJdLFswLDAsIkheM3giXSxbNSwxLCJcXGVxdWl2Il0sWzAsMiwiSHgiXSxbMiwyLCJ4Il0sWzYsMiwiSHgiXSxbMSwyLCJIZiJdLFswLDEsIkheMmYiXSxbMiwzLCJmIiwyXSxbNCw1LCJIZiJdLFs1LDYsImYiXSxbNyw0LCJIXjJmIl0sWzcsOSwiXFxhbHBoYV97SHh9Il0sWzcsOSwiSFxcYWxwaGFfeCIsMix7ImN1cnZlIjo1fV0sWzksMTAsImYiLDJdLFs0LDEwLCJcXGFscGhhX3giLDJdLFsxMCw2LCJcXHZhcmVwc2lsb25feCIsMl0sWzQsOSwiIiwyLHsic2hvcnRlbiI6eyJzb3VyY2UiOjIwLCJ0YXJnZXQiOjIwfSwibGV2ZWwiOjJ9XSxbNSwxMCwiXFxwaGkiLDIseyJzaG9ydGVuIjp7InNvdXJjZSI6MjAsInRhcmdldCI6MjB9LCJsZXZlbCI6Mn1dLFswLDExLCJIXFxhbHBoYV94IiwyXSxbMTEsMiwiXFx2YXJlcHNpbG9uX3tIeH0iLDJdLFs5LDYsImYiLDIseyJjdXJ2ZSI6NX1dLFsxLDExLCJIXFxwaGkiLDAseyJzaG9ydGVuIjp7InNvdXJjZSI6MjAsInRhcmdldCI6MjB9LCJsZXZlbCI6Mn1dLFsxMSwzLCJmIiwyLHsiY3VydmUiOjV9XSxbMTgsMTksIlxcemV0YV94IiwyLHsic2hvcnRlbiI6eyJzb3VyY2UiOjIwLCJ0YXJnZXQiOjIwfX1dLFsxMCwyNywiIiwyLHsic2hvcnRlbiI6eyJzb3VyY2UiOjIwLCJ0YXJnZXQiOjIwfX1dLFsyLDI5LCIiLDIseyJzaG9ydGVuIjp7InNvdXJjZSI6MjAsInRhcmdldCI6MjB9fV1d
		\[\begin{tikzcd}[ampersand replacement=\&,column sep=scriptsize]
			{H^3x} \&\& {H^2x} \&\& Hx \&\& {H^3x} \&\& {H^2x} \&\& \\
			\&\&\&\&\& \equiv \\
			Hx \&\& x \&\& x \&\& Hx \&\& Hx \&\& x
			\arrow["{H^2f}", from=1-1, to=1-3]
			\arrow[""{name=0, anchor=center, inner sep=0}, "{\alpha_{Hx}}", from=1-1, to=3-1]
			\arrow[""{name=1, anchor=center, inner sep=0}, "{H\alpha_x}"', curve={height=30pt}, from=1-1, to=3-1]
			\arrow["Hf", from=1-3, to=1-5]
			\arrow[between={0.2}{0.8}, Rightarrow, nfold, from=1-3, to=3-1]
			\arrow["{\alpha_x}"', from=1-3, to=3-3]
			\arrow["\phi"', between={0.2}{0.8}, Rightarrow, nfold, from=1-5, to=3-3]
			\arrow["f", from=1-5, to=3-5]
			\arrow["{H^2f}", from=1-7, to=1-9]
			\arrow["{H\alpha_x}"', from=1-7, to=3-7]
			\arrow["{H\phi}", between={0.2}{0.8}, Rightarrow, nfold, from=1-9, to=3-7]
			\arrow["Hf", from=1-9, to=3-9]
			\arrow["f"', from=3-1, to=3-3]
			\arrow[""{name=2, anchor=center, inner sep=0}, "f"', curve={height=30pt}, from=3-1, to=3-5]
			\arrow["{\varepsilon_x}"', from=3-3, to=3-5]
			\arrow["{\varepsilon_{Hx}}"', from=3-7, to=3-9]
			\arrow[""{name=3, anchor=center, inner sep=0}, "f"', curve={height=30pt}, from=3-7, to=3-11]
			\arrow["f"', from=3-9, to=3-11]
			\arrow["{\zeta_x}"', between={0.2}{0.8}, Rightarrow, nfold, from=0, to=1]
			\arrow[between={0.2}{0.8}, Rightarrow, nfold, from=3-3, to=2]
			\arrow[between={0.2}{0.8}, Rightarrow, nfold, from=3-9, to=3]
		\end{tikzcd}\cdpunct[8pt]{.}\]
	\end{enumerate}
\end{definition}

\begin{example}
	The functor $\op_2 \colon \nCat{2}_{\leq 2} \rightarrow \nCat{2}_{\leq 2}$ defines a strict $2$-involutive object. Thus, $\op_1 \colon \op_2(\nCat{1}) \rightarrow \nCat{1}$ is a relative strict $1$-involutive object.
\end{example}

One can similarly define the $1$-morphisms and the $2$-morphisms of these relative involutions.

\begin{definition}		\phantom{}
	\begin{enumerate} 
		\item A $1$-morphism $(r,u)$ is given by $r \colon x \rightarrow y$, and a corresponding $2$-cell $u$ that satisfy the following identity:
		% https://q.uiver.app/#q=WzAsMTUsWzEsMCwiSF4yeCJdLFszLDAsIkheMnkiXSxbMiwyLCJIeCJdLFs0LDIsIkh5Il0sWzEsNCwieCJdLFszLDQsInkiXSxbNSwyLCJcXGVxdWl2Il0sWzYsMiwieCJdLFs3LDQsIngiXSxbOCwyLCJ5Il0sWzksNCwieSJdLFsxMCwyLCJIeSJdLFswLDIsIngiXSxbNywwLCJIXjJ4Il0sWzksMCwiSF4yeSJdLFswLDIsIkhmIl0sWzEsMywiSGciXSxbMCwxLCJIXjJyIl0sWzIsMywiSHIiLDJdLFsxLDIsIkh1IiwwLHsic2hvcnRlbiI6eyJzb3VyY2UiOjIwLCJ0YXJnZXQiOjIwfSwibGV2ZWwiOjJ9XSxbMyw1LCJnIl0sWzIsNCwiZiJdLFs0LDUsInIiLDJdLFszLDQsInUiLDAseyJzaG9ydGVuIjp7InNvdXJjZSI6MjAsInRhcmdldCI6MjB9LCJsZXZlbCI6Mn1dLFs4LDEwLCJyIiwyXSxbNyw4LCJcXHZhcmVwc2lsb25feCIsMl0sWzcsOSwiciJdLFs5LDEwLCJcXHZhcmVwc2lsb25feSIsMl0sWzExLDEwLCJnIl0sWzEyLDQsIlxcdmFyZXBzaWxvbl94IiwyXSxbMCwxMiwiXFxhbHBoYV94IiwyXSxbMiwxMiwiXFxwaGlfeCIsMix7InNob3J0ZW4iOnsic291cmNlIjoyMCwidGFyZ2V0IjoyMH0sImxldmVsIjoyfV0sWzE0LDExLCJIZyJdLFsxMSw5LCJcXHBoaV95IiwyLHsic2hvcnRlbiI6eyJzb3VyY2UiOjIwLCJ0YXJnZXQiOjIwfSwibGV2ZWwiOjJ9XSxbMTQsOSwiXFxhbHBoYV95IiwyXSxbMTMsMTQsIkheMnIiXSxbMTMsNywiXFxhbHBoYV94IiwyXSxbOSw4LCIiLDIseyJzaG9ydGVuIjp7InNvdXJjZSI6MjAsInRhcmdldCI6MjB9LCJsZXZlbCI6Mn1dLFsxNCw3LCIiLDAseyJzaG9ydGVuIjp7InNvdXJjZSI6MjAsInRhcmdldCI6MjB9LCJsZXZlbCI6Mn1dXQ==
		\[\begin{tikzcd}[ampersand replacement=\&,cramped,column sep=small]
			\& {H^2x} \&\& {H^2y} \&\&\&\& {H^2x} \&\& {H^2y} \& \\
			\\
			x \&\& Hx \&\& Hy \& \equiv \& x \&\& y \&\& Hy \\
			\\
			\& x \&\& y \&\&\&\& x \&\& y
			\arrow["{H^2r}", from=1-2, to=1-4]
			\arrow["{\alpha_x}"', from=1-2, to=3-1]
			\arrow["Hf", from=1-2, to=3-3]
			\arrow["Hu", between={0.2}{0.8}, Rightarrow, nfold, from=1-4, to=3-3]
			\arrow["Hg", from=1-4, to=3-5]
			\arrow["{H^2r}", from=1-8, to=1-10]
			\arrow["{\alpha_x}"', from=1-8, to=3-7]
			\arrow[between={0.2}{0.8}, Rightarrow, nfold, from=1-10, to=3-7]
			\arrow["{\alpha_y}"', from=1-10, to=3-9]
			\arrow["Hg", from=1-10, to=3-11]
			\arrow["{\varepsilon_x}"', from=3-1, to=5-2]
			\arrow["{\phi_x}"', between={0.2}{0.8}, Rightarrow, nfold, from=3-3, to=3-1]
			\arrow["Hr"', from=3-3, to=3-5]
			\arrow["f", from=3-3, to=5-2]
			\arrow["u", between={0.2}{0.8}, Rightarrow, nfold, from=3-5, to=5-2]
			\arrow["g", from=3-5, to=5-4]
			\arrow["r", from=3-7, to=3-9]
			\arrow["{\varepsilon_x}"', from=3-7, to=5-8]
			\arrow[between={0.2}{0.8}, Rightarrow, nfold, from=3-9, to=5-8]
			\arrow["{\varepsilon_y}"', from=3-9, to=5-10]
			\arrow["{\phi_y}"', between={0.2}{0.8}, Rightarrow, nfold, from=3-11, to=3-9]
			\arrow["g", from=3-11, to=5-10]
			\arrow["r"', from=5-2, to=5-4]
			\arrow["r"', from=5-8, to=5-10]
		\end{tikzcd}\cdpunct[8pt]{.}\]
		\item A $2$-morphism is a $2$-cell $\gamma \colon r \Rightarrow t$ that satisfy the following identity:
		% https://q.uiver.app/#q=WzAsOSxbMCwwLCJIeCJdLFswLDIsIngiXSxbMiwwLCJIeSJdLFsyLDIsInkiXSxbMywxLCJcXGVxdWl2Il0sWzQsMCwiSHgiXSxbNCwyLCJ4Il0sWzYsMCwiSHkiXSxbNiwyLCJ5Il0sWzAsMSwiZiIsMl0sWzEsMywiciJdLFsyLDMsImciXSxbMCwyLCJIciJdLFsyLDEsInUiLDAseyJzaG9ydGVuIjp7InNvdXJjZSI6MjAsInRhcmdldCI6MjB9LCJsZXZlbCI6Mn1dLFs1LDcsIkh0IiwyXSxbNSw2LCJmIiwyXSxbNyw4LCJnIl0sWzYsOCwidCJdLFs3LDYsInYiLDAseyJzaG9ydGVuIjp7InNvdXJjZSI6MjAsInRhcmdldCI6MjB9LCJsZXZlbCI6Mn1dLFsxLDMsInQiLDIseyJjdXJ2ZSI6NX1dLFs1LDcsIkhyIiwwLHsiY3VydmUiOi01fV0sWzEwLDE5LCJcXGdhbW1hIiwwLHsic2hvcnRlbiI6eyJzb3VyY2UiOjIwLCJ0YXJnZXQiOjIwfX1dLFsyMCwxNCwiSFxcZ2FtbWEiLDAseyJzaG9ydGVuIjp7InNvdXJjZSI6MjAsInRhcmdldCI6MjB9fV1d
		\[\begin{tikzcd}[ampersand replacement=\&]
			Hx \&\& Hy \&\& Hx \&\& Hy \\
			\&\&\& \equiv \\
			x \&\& y \&\& x \&\& y
			\arrow["Hr", from=1-1, to=1-3]
			\arrow["f"', from=1-1, to=3-1]
			\arrow["u", between={0.2}{0.8}, Rightarrow, nfold, from=1-3, to=3-1]
			\arrow["g", from=1-3, to=3-3]
			\arrow[""{name=0, anchor=center, inner sep=0}, "Ht"', from=1-5, to=1-7]
			\arrow[""{name=1, anchor=center, inner sep=0}, "Hr", curve={height=-30pt}, from=1-5, to=1-7]
			\arrow["f"', from=1-5, to=3-5]
			\arrow["v", between={0.2}{0.8}, Rightarrow, nfold, from=1-7, to=3-5]
			\arrow["g", from=1-7, to=3-7]
			\arrow[""{name=2, anchor=center, inner sep=0}, "r", from=3-1, to=3-3]
			\arrow[""{name=3, anchor=center, inner sep=0}, "t"', curve={height=30pt}, from=3-1, to=3-3]
			\arrow["t", from=3-5, to=3-7]
			\arrow["{H\gamma}", between={0.2}{0.8}, Rightarrow, nfold, from=1, to=0]
			\arrow["\gamma", between={0.2}{0.8}, Rightarrow, nfold, from=2, to=3]
		\end{tikzcd}\cdpunct[8pt]{.}\]
	\end{enumerate}
\end{definition} 

Therefore, we obtain the bicategory $\nInv{1}(K,H,\alpha,\zeta)$. The axioms can be verified similar as before by using the strictification to encode the desired compositional data.

\begin{proposition}
	The microcosm principle holds for weak $2$-involutive bicategories.
\end{proposition}

Finally, in parallel, we can also write these higher dimensional definitions in the context of inverses. 

\begin{definition}[{\cite[Def.~2.1]{gurski2012biequivintricat}}]
	A weak $2$-adjunction in a tricategory $T$ is
	\begin{enumerate}
		\item two $1$-cell $F \colon K_1 \rightarrow K_2$ and $G \colon K_2 \rightarrow K_1$,
		
		\item two $2$-cells $\epsilon \colon FG \Rightarrow 1_{K_2}$ and $\eta \colon 1_{K_1} \Rightarrow GF$,
		
		\item two $3$-cells $\phi$ and $\psi$ as depicted by the following diagrams:
		% https://q.uiver.app/#q=WzAsMTUsWzEsNCwiS18yIl0sWzAsMiwiS18xIl0sWzMsNCwiS18yIl0sWzIsMiwiS18xIl0sWzUsMiwiS18xIl0sWzgsNCwiS18yIl0sWzQsMywiXFxScmlnaHRhcnJvd157XFxwaGl9Il0sWzIsOCwiRiBcXG90aW1lcyAxX3tLXzF9ICJdLFsyLDksIkYiXSxbNiw4LCIxX3tLXzJ9IFxcb3RpbWVzIEYiXSxbMiw3LCJGIFxcb3RpbWVzKEcgXFxvdGltZXMgRikiXSxbNiw3LCIoRiBcXG90aW1lcyBHKSBcXG90aW1lcyBGIl0sWzMsMCwiRiJdLFswLDYsIkYiXSxbNiw5LCJGIl0sWzEsMCwiRiIsMl0sWzMsMiwiRiJdLFsxLDMsIjFfe0tfMX0iXSxbMCwyLCIxX3tLXzJ9IiwyXSxbMCwzLCJHIiwxXSxbNCw1LCJGIiwxLHsiY3VydmUiOi01fV0sWzcsMTAsIjEgXFxvdGltZXMgXFxldGEiXSxbMTAsMTEsImFeey59Il0sWzExLDksIlxcZXBzaWxvbiBcXG90aW1lcyAxIl0sWzEsMTIsIiIsMSx7InN0eWxlIjp7ImhlYWQiOnsibmFtZSI6Im5vbmUifX19XSxbMTIsMl0sWzEyLDMsIiIsMSx7InNob3J0ZW4iOnsic291cmNlIjoyMCwidGFyZ2V0IjoyMH0sImxldmVsIjoyfV0sWzEsMTMsIiIsMSx7InN0eWxlIjp7ImhlYWQiOnsibmFtZSI6Im5vbmUifX19XSxbMTMsMl0sWzAsMTMsIiIsMSx7InNob3J0ZW4iOnsic291cmNlIjoyMCwidGFyZ2V0IjoyMH0sImxldmVsIjoyfV0sWzQsNSwiRiIsMSx7ImN1cnZlIjo1fV0sWzgsNywicl57Ln0iXSxbOSwxNCwibCJdLFs4LDE0LCIxX0YiLDJdLFsxNywwLCJcXGV0YSIsMSx7InNob3J0ZW4iOnsic291cmNlIjoyMCwidGFyZ2V0IjoyMH19XSxbMywxOCwiXFxlcHNpbG9uIiwxLHsic2hvcnRlbiI6eyJzb3VyY2UiOjIwLCJ0YXJnZXQiOjIwfX1dLFsyMCwzMCwiMV9GIiwxLHsic2hvcnRlbiI6eyJzb3VyY2UiOjIwLCJ0YXJnZXQiOjIwfX1dLFsyMiwzMywiXFxwaGkiLDAseyJzaG9ydGVuIjp7InNvdXJjZSI6MjAsInRhcmdldCI6MjB9fV1d
		\[\begin{tikzcd}[ampersand replacement=\&,column sep=small]
			\&\&\& F \&\&\&\&\& \\
			\\
			{K_1} \&\& {K_1} \&\&\& {K_1} \\
			\&\&\&\& {\Rrightarrow^{\phi}} \\
			\& {K_2} \&\& {K_2} \&\&\&\&\& {K_2} \\
			\\
			F \\
			\&\& {F \otimes(G \otimes F)} \&\&\&\& {(F \otimes G) \otimes F} \\
			\&\& {F \otimes 1_{K_1} } \&\&\&\& {1_{K_2} \otimes F} \\
			\&\& F \&\&\&\& F
			\arrow[between={0.2}{0.8}, Rightarrow, nfold, from=1-4, to=3-3]
			\arrow[from=1-4, to=5-4]
			\arrow[no head, from=3-1, to=1-4]
			\arrow[""{name=0, anchor=center, inner sep=0}, "{1_{K_1}}", from=3-1, to=3-3]
			\arrow["F"', from=3-1, to=5-2]
			\arrow[no head, from=3-1, to=7-1]
			\arrow["F", from=3-3, to=5-4]
			\arrow[""{name=1, anchor=center, inner sep=0}, "F"{description}, curve={height=-30pt}, from=3-6, to=5-9]
			\arrow[""{name=2, anchor=center, inner sep=0}, "F"{description}, curve={height=30pt}, from=3-6, to=5-9]
			\arrow["G"{description}, from=5-2, to=3-3]
			\arrow[""{name=3, anchor=center, inner sep=0}, "{1_{K_2}}"', from=5-2, to=5-4]
			\arrow[between={0.2}{0.8}, Rightarrow, nfold, from=5-2, to=7-1]
			\arrow[from=7-1, to=5-4]
			\arrow[""{name=4, anchor=center, inner sep=0}, "{a^{.}}", from=8-3, to=8-7]
			\arrow["{\epsilon \otimes 1}", from=8-7, to=9-7]
			\arrow["{1 \otimes \eta}", from=9-3, to=8-3]
			\arrow["l", from=9-7, to=10-7]
			\arrow["{r^{.}}", from=10-3, to=9-3]
			\arrow[""{name=5, anchor=center, inner sep=0}, "{1_F}"', from=10-3, to=10-7]
			\arrow["\eta"{description}, between={0.2}{0.8}, Rightarrow, nfold, from=0, to=5-2]
			\arrow["\epsilon"{description}, between={0.2}{0.8}, Rightarrow, nfold, from=3-3, to=3]
			\arrow["{1_F}"{description}, between={0.2}{0.8}, Rightarrow, nfold, from=1, to=2]
			\arrow["\phi", between={0.2}{0.8}, Rightarrow, nfold, from=4, to=5]
		\end{tikzcd}\]
		% https://q.uiver.app/#q=WzAsMTUsWzAsNCwiS18yIl0sWzMsMiwiS18xIl0sWzIsNCwiS18yIl0sWzEsMiwiS18xIl0sWzUsNCwiS18yIl0sWzgsMiwiS18xIl0sWzQsMywiXFxMbGVmdGFycm93XntcXHBzaX0iXSxbNiw4LCJHIFxcb3RpbWVzIDFfe0tfMn0iXSxbMiw5LCJHIl0sWzIsOCwiMV97S18xfSBcXG90aW1lcyBHIl0sWzIsNywiKEcgXFxvdGltZXMgRiApIFxcb3RpbWVzIEciXSxbNiw3LCJHIFxcb3RpbWVzIChGIFxcb3RpbWVzIEcpIl0sWzYsOSwiRyJdLFswLDAsIkciXSxbMyw2LCJHIl0sWzMsMiwiRiIsMV0sWzMsMSwiMV97S18xfSJdLFswLDMsIkciXSxbMiwxLCJHIiwyXSxbNCw1LCJHIiwxLHsiY3VydmUiOi01fV0sWzAsMiwiMV97S18yfSIsMl0sWzksMTAsIlxcZXRhIFxcb3RpbWVzIDEiXSxbMTAsMTEsImEiXSxbMTEsNywiMSBcXG90aW1lcyBcXGVwc2lsb24iXSxbOCw5LCJsXnsufSJdLFs3LDEyLCJyIl0sWzgsMTIsIjFfRyIsMl0sWzAsMTMsIiIsMix7InN0eWxlIjp7ImhlYWQiOnsibmFtZSI6Im5vbmUifX19XSxbMTMsMV0sWzEzLDMsIiIsMix7InNob3J0ZW4iOnsic291cmNlIjoyMCwidGFyZ2V0IjoyMH0sImxldmVsIjoyfV0sWzAsMTQsIiIsMix7InN0eWxlIjp7ImhlYWQiOnsibmFtZSI6Im5vbmUifX19XSxbMTQsMV0sWzIsMTQsIiIsMix7InNob3J0ZW4iOnsic291cmNlIjoyMCwidGFyZ2V0IjoyMH0sImxldmVsIjoyfV0sWzQsNSwiRyIsMSx7ImN1cnZlIjo1fV0sWzE2LDIsIlxcZXRhIiwxLHsic2hvcnRlbiI6eyJzb3VyY2UiOjIwLCJ0YXJnZXQiOjIwfX1dLFszLDIwLCJcXGVwc2lsb24iLDEseyJzaG9ydGVuIjp7InNvdXJjZSI6MjAsInRhcmdldCI6MjB9fV0sWzI2LDIyLCJcXHBzaSIsMix7InNob3J0ZW4iOnsic291cmNlIjoyMCwidGFyZ2V0IjoyMH19XSxbMTksMzMsIjFfRyIsMSx7InNob3J0ZW4iOnsic291cmNlIjoyMCwidGFyZ2V0IjoyMH19XV0=
		\[\begin{tikzcd}[ampersand replacement=\&,column sep=small]
			G \&\&\&\&\&\&\&\& \\
			\\
			\& {K_1} \&\& {K_1} \&\&\&\&\& {K_1} \\
			\&\&\&\& {\Lleftarrow^{\psi}} \\
			{K_2} \&\& {K_2} \&\&\& {K_2} \\
			\\
			\&\&\& G \\
			\&\& {(G \otimes F ) \otimes G} \&\&\&\& {G \otimes (F \otimes G)} \\
			\&\& {1_{K_1} \otimes G} \&\&\&\& {G \otimes 1_{K_2}} \\
			\&\& G \&\&\&\& G
			\arrow[between={0.2}{0.8}, Rightarrow, nfold, from=1-1, to=3-2]
			\arrow[from=1-1, to=3-4]
			\arrow[""{name=0, anchor=center, inner sep=0}, "{1_{K_1}}", from=3-2, to=3-4]
			\arrow["F"{description}, from=3-2, to=5-3]
			\arrow[no head, from=5-1, to=1-1]
			\arrow["G", from=5-1, to=3-2]
			\arrow[""{name=1, anchor=center, inner sep=0}, "{1_{K_2}}"', from=5-1, to=5-3]
			\arrow[no head, from=5-1, to=7-4]
			\arrow["G"', from=5-3, to=3-4]
			\arrow[between={0.2}{0.8}, Rightarrow, nfold, from=5-3, to=7-4]
			\arrow[""{name=2, anchor=center, inner sep=0}, "G"{description}, curve={height=-30pt}, from=5-6, to=3-9]
			\arrow[""{name=3, anchor=center, inner sep=0}, "G"{description}, curve={height=30pt}, from=5-6, to=3-9]
			\arrow[from=7-4, to=3-4]
			\arrow[""{name=4, anchor=center, inner sep=0}, "a", from=8-3, to=8-7]
			\arrow["{1 \otimes \epsilon}", from=8-7, to=9-7]
			\arrow["{\eta \otimes 1}", from=9-3, to=8-3]
			\arrow["r", from=9-7, to=10-7]
			\arrow["{l^{.}}", from=10-3, to=9-3]
			\arrow[""{name=5, anchor=center, inner sep=0}, "{1_G}"', from=10-3, to=10-7]
			\arrow["\eta"{description}, between={0.2}{0.8}, Rightarrow, nfold, from=0, to=5-3]
			\arrow["\epsilon"{description}, between={0.2}{0.8}, Rightarrow, nfold, from=3-2, to=1]
			\arrow["{1_G}"{description}, between={0.2}{0.8}, Rightarrow, nfold, from=2, to=3]
			\arrow["\psi"', between={0.2}{0.8}, Rightarrow, nfold, from=5, to=4]
		\end{tikzcd}\cdpunct[8pt]{,}\]
		\item such that two pasting diagrams as given in \cite[Def.~2.1]{gurski2012biequivintricat} or in \cref{app:weak2invs} hold. 
	\end{enumerate}
\end{definition}

As before, the morphisms of weak $2$-adjunctions can be found in \cref{appsub:weak2invs}.

\begin{definition}
	Let $(F \colon K_1 \rightarrow K_2, G: K_2 \rightarrow K_1, \epsilon,\eta,\phi,\psi)$ be the data of a weak 2-adjunction in the tricategory $\BiCat$. We define a relative (lax) weak $1$-inverse object by
	\begin{enumerate}
		\item two $1$-cells $f \colon Fk_1 \rightarrow k_2 \in K_2$ and $g \colon Gk_2 \rightarrow k_1 \in K_1$,
		
		\item two $2$-cells
		% https://q.uiver.app/#q=WzAsOCxbMiwwLCJGa18xIl0sWzIsMiwia18yIl0sWzAsMCwiRkdrXzIiXSxbMCwyLCJrXzIiXSxbNSwwLCJHa18yIl0sWzUsMiwia18xIl0sWzMsMCwiR0ZrXzEiXSxbMywyLCJrXzEiXSxbMCwxLCJmIl0sWzIsMCwiRihnKSJdLFsyLDMsIlxcZXBzaWxvbl97a18yfSIsMl0sWzMsMSwiXFx2YXJlcHNpbG9uX3trXzJ9IiwyXSxbMCwzLCJcXGFscGhhIiwwLHsic2hvcnRlbiI6eyJzb3VyY2UiOjIwLCJ0YXJnZXQiOjIwfSwibGV2ZWwiOjJ9XSxbNiw0LCJHKGYpIl0sWzQsNSwiZyJdLFs3LDYsIlxcZXRhX3trXzF9Il0sWzcsNSwiXFx2YXJlcHNpbG9uX3trXzF9IiwyXSxbMTYsMTMsIlxcYmV0YSIsMix7InNob3J0ZW4iOnsic291cmNlIjoyMCwidGFyZ2V0IjoyMH19XV0=
		\[\begin{tikzcd}[ampersand replacement=\&]
			{FGk_2} \&\& {Fk_1} \& {GFk_1} \&\& {Gk_2} \\
			\\
			{k_2} \&\& {k_2} \& {k_1} \&\& {k_1}
			\arrow["{F(g)}", from=1-1, to=1-3]
			\arrow["{\epsilon_{k_2}}"', from=1-1, to=3-1]
			\arrow["\alpha", between={0.2}{0.8}, Rightarrow, nfold, from=1-3, to=3-1]
			\arrow["f", from=1-3, to=3-3]
			\arrow[""{name=0, anchor=center, inner sep=0}, "{G(f)}", from=1-4, to=1-6]
			\arrow["g", from=1-6, to=3-6]
			\arrow["{\varepsilon_{k_2}}"', from=3-1, to=3-3]
			\arrow["{\eta_{k_1}}", from=3-4, to=1-4]
			\arrow[""{name=1, anchor=center, inner sep=0}, "{\varepsilon_{k_1}}"', from=3-4, to=3-6]
			\arrow["\beta"', between={0.2}{0.8}, Rightarrow, nfold, from=1, to=0]
		\end{tikzcd}\cdpunct[8pt]{,}\]
		
		\item such that the $2$-cells satisfy the following identities:
		% https://q.uiver.app/#q=WzAsMTMsWzQsMiwiRkdrXzIiXSxbNCw0LCJGIGtfMSJdLFsyLDIsIkZHRmtfMSJdLFswLDQsIkZrXzEiXSxbMCwwLCJGa18xIl0sWzYsMCwia18yIl0sWzYsNCwia18yIl0sWzMsNSwiXFx2ZXF1aXYiXSxbMCw2LCJGa18xIl0sWzYsNiwia18yIl0sWzYsOCwia18yIl0sWzQsOCwiRmtfMSJdLFswLDgsIkZrXzEiXSxbMiwwLCJGRyhmKSIsMV0sWzAsMSwiRmciXSxbMywyLCJGXFxldGFfe2tfMX0iLDFdLFszLDEsIlxcdmFyZXBzaWxvbl97RmtfMX0iLDJdLFsyLDQsIlxcZXBzaWxvbl97RmtfMX0iLDFdLFswLDUsIlxcZXBzaWxvbl97a18yfSIsMV0sWzQsNSwiZiJdLFsxLDYsImYiLDJdLFs1LDYsIlxcdmFyZXBzaWxvbl97a18yfSIsMl0sWzEsNSwiXFxhbHBoYSIsMSx7InNob3J0ZW4iOnsic291cmNlIjoyMCwidGFyZ2V0IjoyMH0sImxldmVsIjoyfV0sWzEyLDgsIlxcdmFyZXBzaWxvbl97RmtfMX0iXSxbOCw5LCJmIl0sWzksMTAsIlxcdmFyZXBzaWxvbl97a18yfSJdLFsxMSwxMCwiZiIsMl0sWzEyLDExLCJcXHZhcmVwc2lsb25fe0ZrXzF9IiwyXSxbOCwxMCwiZiIsMV0sWzgsMTEsIlxcdmFyZXBzaWxvbl97RmtfMX0iLDFdLFszLDQsIlxcdmFyZXBzaWxvbl97RmtfMX0iXSxbMTYsMTMsIkZcXGJldGEiLDEseyJzaG9ydGVuIjp7InNvdXJjZSI6MjAsInRhcmdldCI6MjB9fV0sWzEzLDE5LCIiLDAseyJzaG9ydGVuIjp7InNvdXJjZSI6MjAsInRhcmdldCI6MjB9fV0sWzIsMzAsIlxccGhpIiwxLHsic2hvcnRlbiI6eyJzb3VyY2UiOjIwLCJ0YXJnZXQiOjIwfX1dXQ==
		\[\begin{tikzcd}[ampersand replacement=\&,cramped]
			{Fk_1} \&\&\&\&\&\& {k_2} \\
			\\
			\&\& {FGFk_1} \&\& {FGk_2} \\
			\\
			{Fk_1} \&\&\&\& {F k_1} \&\& {k_2} \\
			\&\&\& \vequiv \\
			{Fk_1} \&\&\&\&\&\& {k_2} \\
			\\
			{Fk_1} \&\&\&\& {Fk_1} \&\& {k_2}
			\arrow[""{name=0, anchor=center, inner sep=0}, "f", from=1-1, to=1-7]
			\arrow["{\varepsilon_{k_2}}"', from=1-7, to=5-7]
			\arrow["{\epsilon_{Fk_1}}"{description}, from=3-3, to=1-1]
			\arrow[""{name=1, anchor=center, inner sep=0}, "{FG(f)}"{description}, from=3-3, to=3-5]
			\arrow["{\epsilon_{k_2}}"{description}, from=3-5, to=1-7]
			\arrow["Fg", from=3-5, to=5-5]
			\arrow[""{name=2, anchor=center, inner sep=0}, "{\varepsilon_{Fk_1}}", from=5-1, to=1-1]
			\arrow["{F\eta_{k_1}}"{description}, from=5-1, to=3-3]
			\arrow[""{name=3, anchor=center, inner sep=0}, "{\varepsilon_{Fk_1}}"', from=5-1, to=5-5]
			\arrow["\alpha"{description}, between={0.2}{0.8}, Rightarrow, nfold, from=5-5, to=1-7]
			\arrow["f"', from=5-5, to=5-7]
			\arrow["f", from=7-1, to=7-7]
			\arrow["{\varepsilon_{Fk_1}}"{description}, from=7-1, to=9-5]
			\arrow["f"{description}, from=7-1, to=9-7]
			\arrow["{\varepsilon_{k_2}}", from=7-7, to=9-7]
			\arrow["{\varepsilon_{Fk_1}}", from=9-1, to=7-1]
			\arrow["{\varepsilon_{Fk_1}}"', from=9-1, to=9-5]
			\arrow["f"', from=9-5, to=9-7]
			\arrow["\phi"{description}, between={0.2}{0.8}, Rightarrow, nfold, from=3-3, to=2]
			\arrow[between={0.2}{0.8}, Rightarrow, nfold, from=1, to=0]
			\arrow["{F\beta}"{description}, between={0.2}{0.8}, Rightarrow, nfold, from=3, to=1]
		\end{tikzcd}\]
		% https://q.uiver.app/#q=WzAsMTQsWzUsMiwiR0ZrXzEiXSxbNSw0LCJHa18yIl0sWzEsMiwiR0ZHa18yIl0sWzEsNCwiR2tfMiJdLFs3LDQsImtfMSJdLFs3LDAsImtfMSJdLFsxLDAsIkdrXzIiXSxbNyw2LCJrXzEiXSxbNyw4LCJrXzEiXSxbMSw2LCJHa18yIl0sWzEsOCwiR2tfMiJdLFs0LDgsIkdrXzIiXSxbNCw1LCJcXHZlcXVpdiJdLFswLDcsIkdGR2tfMiJdLFswLDEsIkdmIiwxXSxbMiwwLCJHRihnKSIsMV0sWzIsMywiR1xcZXBzaWxvbl97a18yfSIsMl0sWzMsMSwiXFx2YXJlcHNpbG9uX3tHa18yfSIsMl0sWzAsMywiR1xcYWxwaGEiLDEseyJzaG9ydGVuIjp7InNvdXJjZSI6MjAsInRhcmdldCI6MjB9LCJsZXZlbCI6Mn1dLFsxLDQsImciLDJdLFs1LDAsIlxcZXRhX3trXzF9IiwxXSxbNSw0LCJcXHZhcmVwc2lsb25fe2tfMX0iXSxbNiw1LCJnIl0sWzYsMiwiXFxldGFfe0drXzJ9IiwyXSxbOSw4LCJnIiwxXSxbNyw4LCJcXHZhcmVwc2lsb25fe2tfMX0iLDFdLFs5LDcsImciXSxbMTEsOCwiZyIsMl0sWzksMTAsIlxcdmFyZXBzaWxvbl97R2tfMn0iLDFdLFsxMCwxMSwiXFx2YXJlcHNpbG9uX3tHa18yfSIsMl0sWzksMTEsIlxcdmFyZXBzaWxvbl97R2tfMn0iLDFdLFs5LDEzLCJcXGV0YV97R2tfMn0iLDJdLFsxMywxMCwiR1xcZXBzaWxvbl97a18yfSIsMl0sWzIyLDE1LCIiLDAseyJzaG9ydGVuIjp7InNvdXJjZSI6MjAsInRhcmdldCI6MjB9fV0sWzIxLDE0LCJcXGJldGEiLDEseyJzaG9ydGVuIjp7InNvdXJjZSI6MjAsInRhcmdldCI6MjB9fV0sWzI4LDEzLCJcXHBzaSIsMix7InNob3J0ZW4iOnsic291cmNlIjoyMCwidGFyZ2V0IjoyMH19XV0=
		\[\begin{tikzcd}[ampersand replacement=\&,cramped]
			\& {Gk_2} \&\&\&\&\&\& {k_1} \\
			\\
			\& {GFGk_2} \&\&\&\& {GFk_1} \\
			\\
			\& {Gk_2} \&\&\&\& {Gk_2} \&\& {k_1} \\
			\&\&\&\& \vequiv \\
			\& {Gk_2} \&\&\&\&\&\& {k_1} \\
			{GFGk_2} \\
			\& {Gk_2} \&\&\& {Gk_2} \&\&\& {k_1}
			\arrow[""{name=0, anchor=center, inner sep=0}, "g", from=1-2, to=1-8]
			\arrow["{\eta_{Gk_2}}"', from=1-2, to=3-2]
			\arrow["{\eta_{k_1}}"{description}, from=1-8, to=3-6]
			\arrow[""{name=1, anchor=center, inner sep=0}, "{\varepsilon_{k_1}}", from=1-8, to=5-8]
			\arrow[""{name=2, anchor=center, inner sep=0}, "{GF(g)}"{description}, from=3-2, to=3-6]
			\arrow["{G\epsilon_{k_2}}"', from=3-2, to=5-2]
			\arrow["{G\alpha}"{description}, between={0.2}{0.8}, Rightarrow, nfold, from=3-6, to=5-2]
			\arrow[""{name=3, anchor=center, inner sep=0}, "Gf"{description}, from=3-6, to=5-6]
			\arrow["{\varepsilon_{Gk_2}}"', from=5-2, to=5-6]
			\arrow["g"', from=5-6, to=5-8]
			\arrow["g", from=7-2, to=7-8]
			\arrow["{\eta_{Gk_2}}"', from=7-2, to=8-1]
			\arrow[""{name=4, anchor=center, inner sep=0}, "{\varepsilon_{Gk_2}}"{description}, from=7-2, to=9-2]
			\arrow["{\varepsilon_{Gk_2}}"{description}, from=7-2, to=9-5]
			\arrow["g"{description}, from=7-2, to=9-8]
			\arrow["{\varepsilon_{k_1}}"{description}, from=7-8, to=9-8]
			\arrow["{G\epsilon_{k_2}}"', from=8-1, to=9-2]
			\arrow["{\varepsilon_{Gk_2}}"', from=9-2, to=9-5]
			\arrow["g"', from=9-5, to=9-8]
			\arrow[between={0.2}{0.8}, Rightarrow, nfold, from=0, to=2]
			\arrow["\beta"{description}, between={0.2}{0.8}, Rightarrow, nfold, from=1, to=3]
			\arrow["\psi"', between={0.2}{0.8}, Rightarrow, nfold, from=4, to=8-1]
		\end{tikzcd}\cdpunct[8pt]{.}\]
	\end{enumerate}
\end{definition}

\begin{definition}
	Let $(f\colon Fk_1 \rightarrow k_2, g \colon Gk_2 \rightarrow k_1, \alpha,\beta)$ and $(f'\colon Fk'_1 \rightarrow k'_2, g' \colon Gk'_2 \rightarrow k'_1, \alpha',\beta')$ be two relative inverse objects. The data of a (lax) $1$-morphism is given by
	\begin{enumerate}
		\item two $1$-cells $k_1 \rightarrow k'_1$ and $k_2 \rightarrow k'_2$,
		
		\item two $2$-cells
		% https://q.uiver.app/#q=WzAsOCxbMCwwLCJGa18xIl0sWzIsMCwiRmsnXzEiXSxbMCwyLCJrXzIiXSxbMiwyLCJrJ18yIl0sWzQsMCwiR2tfMiJdLFs2LDAsIkdrJ18yIl0sWzQsMiwia18xIl0sWzYsMiwiaydfMSJdLFswLDFdLFswLDIsImYiLDJdLFsxLDMsImYnIl0sWzIsM10sWzEsMiwiIiwwLHsic2hvcnRlbiI6eyJzb3VyY2UiOjIwLCJ0YXJnZXQiOjIwfSwibGV2ZWwiOjJ9XSxbNSw3LCJnIl0sWzQsNiwiZyIsMl0sWzQsNV0sWzYsN10sWzUsNiwiIiwwLHsic2hvcnRlbiI6eyJzb3VyY2UiOjIwLCJ0YXJnZXQiOjIwfSwibGV2ZWwiOjJ9XV0=
		\[\begin{tikzcd}[ampersand replacement=\&]
			{Fk_1} \&\& {Fk'_1} \&\& {Gk_2} \&\& {Gk'_2} \\
			\\
			{k_2} \&\& {k'_2} \&\& {k_1} \&\& {k'_1}
			\arrow[from=1-1, to=1-3]
			\arrow["f"', from=1-1, to=3-1]
			\arrow[between={0.2}{0.8}, Rightarrow, nfold, from=1-3, to=3-1]
			\arrow["{f'}", from=1-3, to=3-3]
			\arrow[from=1-5, to=1-7]
			\arrow["g"', from=1-5, to=3-5]
			\arrow[between={0.2}{0.8}, Rightarrow, nfold, from=1-7, to=3-5]
			\arrow["g", from=1-7, to=3-7]
			\arrow[from=3-1, to=3-3]
			\arrow[from=3-5, to=3-7]
		\end{tikzcd}\cdpunct[8pt]{,}\]
		
		\item such that the following two identities hold:
		% https://q.uiver.app/#q=WzAsMTUsWzAsMiwiR0ZrXzEiXSxbMiwyLCJHRmsnXzEiXSxbMCw0LCJHa18yIl0sWzIsNCwiR2snXzIiXSxbMSw2LCJrXzEiXSxbMyw2LCJrJ18xIl0sWzEsMCwia18xIl0sWzMsMCwiaydfMSJdLFs0LDMsIlxcZXF1aXYiXSxbNiwwLCJrXzEiXSxbNSwyLCJHRmtfMSJdLFs1LDQsIkdrXzIiXSxbNiw2LCJrXzEiXSxbOCwwLCJrXzEnIl0sWzgsNiwiaydfMSJdLFswLDFdLFswLDIsIkdmIiwyXSxbMSwzLCJHZiciLDFdLFsyLDNdLFsxLDIsIiIsMCx7InNob3J0ZW4iOnsic291cmNlIjoyMCwidGFyZ2V0IjoyMH0sImxldmVsIjoyfV0sWzMsNSwiZyciLDFdLFsyLDQsImciLDJdLFs0LDVdLFszLDQsIiIsMCx7InNob3J0ZW4iOnsic291cmNlIjoyMCwidGFyZ2V0IjoyMH0sImxldmVsIjoyfV0sWzcsMSwiXFxldGFfe2snXzF9IiwxXSxbNiwwLCJcXGV0YV97a18xfSIsMl0sWzYsN10sWzcsMCwiIiwwLHsic2hvcnRlbiI6eyJzb3VyY2UiOjIwLCJ0YXJnZXQiOjIwfSwibGV2ZWwiOjJ9XSxbNyw1LCJcXHZhcmVwc2lsb25fe2snXzF9Il0sWzksMTIsIlxcdmFyZXBzaWxvbl97a18xfSJdLFs5LDEwLCJcXGV0YV97a18xfSIsMl0sWzEwLDExLCJHZiIsMl0sWzExLDEyLCJnIiwyXSxbOSwxM10sWzEyLDE0XSxbMTMsMTQsIlxcdmFyZXBzaWxvbl97aydfMX0iXSxbOSwxNF0sWzI4LDE3LCJcXGJldGEnIiwxLHsic2hvcnRlbiI6eyJzb3VyY2UiOjIwLCJ0YXJnZXQiOjIwfX1dLFsyOSwzMSwiXFxiZXRhIiwxLHsic2hvcnRlbiI6eyJzb3VyY2UiOjIwLCJ0YXJnZXQiOjIwfX1dXQ==
		\[\begin{tikzcd}[ampersand replacement=\&,column sep=scriptsize]
			\& {k_1} \&\& {k'_1} \&\&\& {k_1} \&\& {k_1'} \\
			\\
			{GFk_1} \&\& {GFk'_1} \&\&\& {GFk_1} \\
			\&\&\&\& \equiv \\
			{Gk_2} \&\& {Gk'_2} \&\&\& {Gk_2} \\
			\\
			\& {k_1} \&\& {k'_1} \&\&\& {k_1} \&\& {k'_1}
			\arrow[from=1-2, to=1-4]
			\arrow["{\eta_{k_1}}"', from=1-2, to=3-1]
			\arrow[between={0.2}{0.8}, Rightarrow, nfold, from=1-4, to=3-1]
			\arrow["{\eta_{k'_1}}"{description}, from=1-4, to=3-3]
			\arrow[""{name=0, anchor=center, inner sep=0}, "{\varepsilon_{k'_1}}", from=1-4, to=7-4]
			\arrow[from=1-7, to=1-9]
			\arrow["{\eta_{k_1}}"', from=1-7, to=3-6]
			\arrow[""{name=1, anchor=center, inner sep=0}, "{\varepsilon_{k_1}}", from=1-7, to=7-7]
			\arrow[from=1-7, to=7-9]
			\arrow["{\varepsilon_{k'_1}}", from=1-9, to=7-9]
			\arrow[from=3-1, to=3-3]
			\arrow["Gf"', from=3-1, to=5-1]
			\arrow[between={0.2}{0.8}, Rightarrow, nfold, from=3-3, to=5-1]
			\arrow[""{name=2, anchor=center, inner sep=0}, "{Gf'}"{description}, from=3-3, to=5-3]
			\arrow[""{name=3, anchor=center, inner sep=0}, "Gf"', from=3-6, to=5-6]
			\arrow[from=5-1, to=5-3]
			\arrow["g"', from=5-1, to=7-2]
			\arrow[between={0.2}{0.8}, Rightarrow, nfold, from=5-3, to=7-2]
			\arrow["{g'}"{description}, from=5-3, to=7-4]
			\arrow["g"', from=5-6, to=7-7]
			\arrow[from=7-2, to=7-4]
			\arrow[from=7-7, to=7-9]
			\arrow["{\beta'}"{description}, between={0.2}{0.8}, Rightarrow, nfold, from=0, to=2]
			\arrow["\beta"{description}, between={0.2}{0.8}, Rightarrow, nfold, from=1, to=3]
		\end{tikzcd}\]
		
		% https://q.uiver.app/#q=WzAsMTUsWzEsNCwia18yIl0sWzMsNCwiaydfMiJdLFsyLDIsIkZrXzEiXSxbNCwyLCJGaydfMSJdLFsxLDAsIkZHa18yIl0sWzMsMCwiRkdrJ18yIl0sWzAsMiwia18yIl0sWzUsMiwiXFxlcXVpdiJdLFs2LDIsImtfMiJdLFs3LDAsIkZHa18yIl0sWzksMCwiRkdrJ18yIl0sWzEwLDIsIkZrJ18xIl0sWzcsNCwia18yIl0sWzksNCwiaydfMiJdLFs4LDIsImsnXzIiXSxbNCw2LCJcXGVwc2lsb25fe2tfMn0iLDJdLFs0LDVdLFsyLDAsImYiLDFdLFszLDEsImYnIl0sWzUsMywiRmcnIl0sWzQsMiwiRmciLDFdLFswLDFdLFsyLDNdLFszLDAsIiIsMCx7InNob3J0ZW4iOnsic291cmNlIjoyMCwidGFyZ2V0IjoyMH0sImxldmVsIjoyfV0sWzUsMiwiIiwwLHsic2hvcnRlbiI6eyJzb3VyY2UiOjIwLCJ0YXJnZXQiOjIwfSwibGV2ZWwiOjJ9XSxbNiwwLCJcXHZhcmVwc2lsb25fe2tfMn0iLDJdLFsyLDYsIlxcYWxwaGEiLDEseyJzaG9ydGVuIjp7InNvdXJjZSI6MjAsInRhcmdldCI6MjB9LCJsZXZlbCI6Mn1dLFsxMCwxNCwiXFxlcHNpbG9uX3trXzF9IiwxXSxbMTAsMTEsIkZnJyJdLFsxNCwxMywiXFx2YXJlcHNpbG9uX3trJ18yfSIsMV0sWzExLDEzLCJmJyJdLFs5LDEwXSxbOCwxNF0sWzksOCwiXFxlcHNpbG9uX3trXzF9IiwyXSxbOCwxMiwiXFx2YXJlcHNpbG9uX3trXzJ9IiwyXSxbMTIsMTNdLFsxMSwxNCwiXFxhbHBoYSciLDEseyJzaG9ydGVuIjp7InNvdXJjZSI6MjAsInRhcmdldCI6MjB9LCJsZXZlbCI6Mn1dLFsxNCwxMiwiIiwxLHsic2hvcnRlbiI6eyJzb3VyY2UiOjIwLCJ0YXJnZXQiOjIwfSwibGV2ZWwiOjJ9XSxbMTAsOCwiIiwxLHsic2hvcnRlbiI6eyJzb3VyY2UiOjIwLCJ0YXJnZXQiOjIwfSwibGV2ZWwiOjJ9XV0=
		\[\begin{tikzcd}[ampersand replacement=\&,column sep=tiny]
			\& {FGk_2} \&\& {FGk'_2} \&\&\&\& {FGk_2} \&\& {FGk'_2} \& \\
			\\
			{k_2} \&\& {Fk_1} \&\& {Fk'_1} \& \equiv \& {k_2} \&\& {k'_2} \&\& {Fk'_1} \\
			\\
			\& {k_2} \&\& {k'_2} \&\&\&\& {k_2} \&\& {k'_2}
			\arrow[from=1-2, to=1-4]
			\arrow["{\epsilon_{k_2}}"', from=1-2, to=3-1]
			\arrow["Fg"{description}, from=1-2, to=3-3]
			\arrow[between={0.2}{0.8}, Rightarrow, nfold, from=1-4, to=3-3]
			\arrow["{Fg'}", from=1-4, to=3-5]
			\arrow[from=1-8, to=1-10]
			\arrow["{\epsilon_{k_1}}"', from=1-8, to=3-7]
			\arrow[between={0.2}{0.8}, Rightarrow, nfold, from=1-10, to=3-7]
			\arrow["{\epsilon_{k_1}}"{description}, from=1-10, to=3-9]
			\arrow["{Fg'}", from=1-10, to=3-11]
			\arrow["{\varepsilon_{k_2}}"', from=3-1, to=5-2]
			\arrow["\alpha"{description}, between={0.2}{0.8}, Rightarrow, nfold, from=3-3, to=3-1]
			\arrow[from=3-3, to=3-5]
			\arrow["f"{description}, from=3-3, to=5-2]
			\arrow[between={0.2}{0.8}, Rightarrow, nfold, from=3-5, to=5-2]
			\arrow["{f'}", from=3-5, to=5-4]
			\arrow[from=3-7, to=3-9]
			\arrow["{\varepsilon_{k_2}}"', from=3-7, to=5-8]
			\arrow[between={0.2}{0.8}, Rightarrow, nfold, from=3-9, to=5-8]
			\arrow["{\varepsilon_{k'_2}}"{description}, from=3-9, to=5-10]
			\arrow["{\alpha'}"{description}, between={0.2}{0.8}, Rightarrow, nfold, from=3-11, to=3-9]
			\arrow["{f'}", from=3-11, to=5-10]
			\arrow[from=5-2, to=5-4]
			\arrow[from=5-8, to=5-10]
		\end{tikzcd}\cdpunct[8pt]{.}\]
		
	\end{enumerate}
\end{definition}

\begin{definition}
	The data of a $2$-morphism between two $1$-morphisms is given
	\begin{enumerate}
		\item two $2$-cells
		% https://q.uiver.app/#q=WzAsNCxbMCwwLCJrXzEiXSxbMiwwLCJrJ18xIl0sWzQsMCwia18yIl0sWzYsMCwiaydfMiJdLFswLDEsIiIsMCx7ImN1cnZlIjotMn1dLFswLDEsIiIsMix7ImN1cnZlIjoyfV0sWzIsMywiIiwyLHsiY3VydmUiOi0yfV0sWzIsMywiIiwwLHsiY3VydmUiOjJ9XSxbNCw1LCIiLDAseyJzaG9ydGVuIjp7InNvdXJjZSI6MjAsInRhcmdldCI6MjB9fV0sWzYsNywiIiwyLHsic2hvcnRlbiI6eyJzb3VyY2UiOjIwLCJ0YXJnZXQiOjIwfX1dXQ==
		\[\begin{tikzcd}[ampersand replacement=\&]
			{k_1} \&\& {k'_1} \&\& {k_2} \&\& {k'_2}
			\arrow[""{name=0, anchor=center, inner sep=0}, curve={height=-12pt}, from=1-1, to=1-3]
			\arrow[""{name=1, anchor=center, inner sep=0}, curve={height=12pt}, from=1-1, to=1-3]
			\arrow[""{name=2, anchor=center, inner sep=0}, curve={height=-12pt}, from=1-5, to=1-7]
			\arrow[""{name=3, anchor=center, inner sep=0}, curve={height=12pt}, from=1-5, to=1-7]
			\arrow[between={0.2}{0.8}, Rightarrow, nfold, from=0, to=1]
			\arrow[between={0.2}{0.8}, Rightarrow, nfold, from=2, to=3]
		\end{tikzcd}\cdpunct[2pt]{,}\]
		\item such that the following two identities hold:
		% https://q.uiver.app/#q=WzAsMTgsWzAsMCwiRmtfMSJdLFsyLDAsIkZrJ18xIl0sWzQsMiwia18yIl0sWzYsMiwiaydfMiJdLFs0LDAsIkZrXzEiXSxbNiwwLCJGaydfMSJdLFswLDIsImtfMiJdLFsyLDIsImsnXzIiXSxbMywxLCJcXGVxdWl2Il0sWzAsNCwiR2tfMiJdLFsyLDQsIkdrJ18yIl0sWzAsNiwia18xIl0sWzIsNiwiaydfMSJdLFs0LDQsIkdrXzIiXSxbNCw2LCJrXzEiXSxbNiw0LCJHaydfMiJdLFs2LDYsImsnXzEiXSxbMyw1LCJcXGVxdWl2Il0sWzAsMSwiIiwwLHsiY3VydmUiOi00fV0sWzAsMV0sWzIsM10sWzIsMywiIiwwLHsiY3VydmUiOjR9XSxbNCw1XSxbNCwyLCJmIiwyXSxbNSwzLCJmJyJdLFs1LDIsIiIsMix7InNob3J0ZW4iOnsic291cmNlIjoyMCwidGFyZ2V0IjoyMH0sImxldmVsIjoyfV0sWzAsNiwiZiIsMl0sWzEsNywiZiciXSxbNiw3XSxbMSw2LCIiLDIseyJzaG9ydGVuIjp7InNvdXJjZSI6MjAsInRhcmdldCI6MjB9LCJsZXZlbCI6Mn1dLFs5LDExLCJnIiwyXSxbMTEsMTJdLFs5LDEwXSxbMTAsMTIsImcnIl0sWzE1LDE2LCJnJyJdLFsxMywxNCwiZyIsMl0sWzE0LDE2XSxbMTMsMTVdLFs5LDEwLCIiLDEseyJjdXJ2ZSI6LTR9XSxbMTQsMTYsIiIsMSx7ImN1cnZlIjo0fV0sWzEwLDExLCIiLDAseyJzaG9ydGVuIjp7InNvdXJjZSI6MjAsInRhcmdldCI6MjB9LCJsZXZlbCI6Mn1dLFsxNSwxNCwiIiwxLHsic2hvcnRlbiI6eyJzb3VyY2UiOjIwLCJ0YXJnZXQiOjIwfSwibGV2ZWwiOjJ9XSxbMTgsMTksIiIsMCx7InNob3J0ZW4iOnsic291cmNlIjoyMCwidGFyZ2V0IjoyMH19XSxbMjAsMjEsIiIsMix7InNob3J0ZW4iOnsic291cmNlIjoyMCwidGFyZ2V0IjoyMH19XSxbMzgsMzIsIiIsMSx7InNob3J0ZW4iOnsic291cmNlIjoyMCwidGFyZ2V0IjoyMH19XSxbMzYsMzksIiIsMSx7InNob3J0ZW4iOnsic291cmNlIjoyMCwidGFyZ2V0IjoyMH19XV0=
		\[\begin{tikzcd}[ampersand replacement=\&]
			{Fk_1} \&\& {Fk'_1} \&\& {Fk_1} \&\& {Fk'_1} \\
			\&\&\& \equiv \\
			{k_2} \&\& {k'_2} \&\& {k_2} \&\& {k'_2} \\
			\\
			{Gk_2} \&\& {Gk'_2} \&\& {Gk_2} \&\& {Gk'_2} \\
			\&\&\& \equiv \\
			{k_1} \&\& {k'_1} \&\& {k_1} \&\& {k'_1}
			\arrow[""{name=0, anchor=center, inner sep=0}, curve={height=-24pt}, from=1-1, to=1-3]
			\arrow[""{name=1, anchor=center, inner sep=0}, from=1-1, to=1-3]
			\arrow["f"', from=1-1, to=3-1]
			\arrow[between={0.2}{0.8}, Rightarrow, nfold, from=1-3, to=3-1]
			\arrow["{f'}", from=1-3, to=3-3]
			\arrow[from=1-5, to=1-7]
			\arrow["f"', from=1-5, to=3-5]
			\arrow[between={0.2}{0.8}, Rightarrow, nfold, from=1-7, to=3-5]
			\arrow["{f'}", from=1-7, to=3-7]
			\arrow[from=3-1, to=3-3]
			\arrow[""{name=2, anchor=center, inner sep=0}, from=3-5, to=3-7]
			\arrow[""{name=3, anchor=center, inner sep=0}, curve={height=24pt}, from=3-5, to=3-7]
			\arrow[""{name=4, anchor=center, inner sep=0}, from=5-1, to=5-3]
			\arrow[""{name=5, anchor=center, inner sep=0}, curve={height=-24pt}, from=5-1, to=5-3]
			\arrow["g"', from=5-1, to=7-1]
			\arrow[between={0.2}{0.8}, Rightarrow, nfold, from=5-3, to=7-1]
			\arrow["{g'}", from=5-3, to=7-3]
			\arrow[from=5-5, to=5-7]
			\arrow["g"', from=5-5, to=7-5]
			\arrow[between={0.2}{0.8}, Rightarrow, nfold, from=5-7, to=7-5]
			\arrow["{g'}", from=5-7, to=7-7]
			\arrow[from=7-1, to=7-3]
			\arrow[""{name=6, anchor=center, inner sep=0}, from=7-5, to=7-7]
			\arrow[""{name=7, anchor=center, inner sep=0}, curve={height=24pt}, from=7-5, to=7-7]
			\arrow[between={0.2}{0.8}, Rightarrow, nfold, from=0, to=1]
			\arrow[between={0.2}{0.8}, Rightarrow, nfold, from=2, to=3]
			\arrow[between={0.2}{0.8}, Rightarrow, nfold, from=5, to=4]
			\arrow[between={0.2}{0.8}, Rightarrow, nfold, from=6, to=7]
		\end{tikzcd}\cdpunct[8pt]{.}\]
	\end{enumerate}
\end{definition}

\begin{proposition}
	The microcosm principle holds for weak $2$-adjunction of bicategories.
\end{proposition}

We have now developed all the desired definitions to fully express the dualities of monoidal structures, albiet at a greater level of generality than necessary. For all our purposes, we shall only require $T = \Gray$. 

As a closing remark, we note that so far we have only employed the microcosm principle in doubly nested structures. This is in accordance with how the original statement is written and demonstrated. However, there are instances where we can achieve triply/higher nested structures.

\begin{example}
	The dual vector space $1$-functor $(-)^{*} \colon \op_{1}\left(\mathbf{Vect}_{k}^{f.d.}\right) \rightarrow \mathbf{Vect}_{k}^{f.d.}$ defines an involutive object relative to the strict involutive object $\op_1 \colon \op_2(\nCat{1}) \rightarrow \nCat{1}$. The structure cells are simply $\eta_V \colon V \rightarrow V^{**} \in \mathbf{Vect}_{k}^{f.d}$ as the unit and its opposite cell $\op_1(\eta_V)$ as the counit. One could view $(-)^{*}$ as a contravariant autoequivalence on $\mathbf{Vect}_k^{f.d.}$.
\end{example}

In this example, we have a top ($2$-truncated) $3$-dimensional involution $\op_2 \colon \nCat{2}_{\leq 2} \rightarrow \nCat{2}_{\leq 2}$, inside which we define a ($1$-dimensional promoted to) $2$-dimensional involution $\op_1$, and finally inside which we define a $1$-dimensional involution $(-)^*$. In general, it is not possible to obtain such triply (or higher) nested structures just by simply considering the microcosm principle. 